% concatenated from main.tex, math_commands.tex
\documentclass{article} % For LaTeX2e
\PassOptionsToPackage{table}{xcolor}
\usepackage{iclr2025_conference,times}
% Optional math commands from https://github.com/goodfeli/dlbook_notation.
%%%%% NEW MATH DEFINITIONS %%%%%

\usepackage{amsmath,amsfonts,bm}

% Mark sections of captions for referring to divisions of figures
\newcommand{\figleft}{{\em (Left)}}
\newcommand{\figcenter}{{\em (Center)}}
\newcommand{\figright}{{\em (Right)}}
\newcommand{\figtop}{{\em (Top)}}
\newcommand{\figbottom}{{\em (Bottom)}}
\newcommand{\captiona}{{\em (a)}}
\newcommand{\captionb}{{\em (b)}}
\newcommand{\captionc}{{\em (c)}}
\newcommand{\captiond}{{\em (d)}}

% Highlight a newly defined term
\newcommand{\newterm}[1]{{\bf #1}}


% Figure reference, lower-case.
\def\figref#1{figure~\ref{#1}}
% Figure reference, capital. For start of sentence
\def\Figref#1{Figure~\ref{#1}}
\def\twofigref#1#2{figures \ref{#1} and \ref{#2}}
\def\quadfigref#1#2#3#4{figures \ref{#1}, \ref{#2}, \ref{#3} and \ref{#4}}
% Section reference, lower-case.
\def\secref#1{section~\ref{#1}}
% Section reference, capital.
\def\Secref#1{Section~\ref{#1}}
% Reference to two sections.
\def\twosecrefs#1#2{sections \ref{#1} and \ref{#2}}
% Reference to three sections.
\def\secrefs#1#2#3{sections \ref{#1}, \ref{#2} and \ref{#3}}
% Reference to an equation, lower-case.
\def\eqref#1{equation~\ref{#1}}
% Reference to an equation, upper case
\def\Eqref#1{Equation~\ref{#1}}
% A raw reference to an equation---avoid using if possible
\def\plaineqref#1{\ref{#1}}
% Reference to a chapter, lower-case.
\def\chapref#1{chapter~\ref{#1}}
% Reference to an equation, upper case.
\def\Chapref#1{Chapter~\ref{#1}}
% Reference to a range of chapters
\def\rangechapref#1#2{chapters\ref{#1}--\ref{#2}}
% Reference to an algorithm, lower-case.
\def\algref#1{algorithm~\ref{#1}}
% Reference to an algorithm, upper case.
\def\Algref#1{Algorithm~\ref{#1}}
\def\twoalgref#1#2{algorithms \ref{#1} and \ref{#2}}
\def\Twoalgref#1#2{Algorithms \ref{#1} and \ref{#2}}
% Reference to a part, lower case
\def\partref#1{part~\ref{#1}}
% Reference to a part, upper case
\def\Partref#1{Part~\ref{#1}}
\def\twopartref#1#2{parts \ref{#1} and \ref{#2}}

\def\ceil#1{\lceil #1 \rceil}
\def\floor#1{\lfloor #1 \rfloor}
\def\1{\bm{1}}
\newcommand{\train}{\mathcal{D}}
\newcommand{\valid}{\mathcal{D_{\mathrm{valid}}}}
\newcommand{\test}{\mathcal{D_{\mathrm{test}}}}

\def\eps{{\epsilon}}


% Random variables
\def\reta{{\textnormal{$\eta$}}}
\def\ra{{\textnormal{a}}}
\def\rb{{\textnormal{b}}}
\def\rc{{\textnormal{c}}}
\def\rd{{\textnormal{d}}}
\def\re{{\textnormal{e}}}
\def\rf{{\textnormal{f}}}
\def\rg{{\textnormal{g}}}
\def\rh{{\textnormal{h}}}
\def\ri{{\textnormal{i}}}
\def\rj{{\textnormal{j}}}
\def\rk{{\textnormal{k}}}
\def\rl{{\textnormal{l}}}
% rm is already a command, just don't name any random variables m
\def\rn{{\textnormal{n}}}
\def\ro{{\textnormal{o}}}
\def\rp{{\textnormal{p}}}
\def\rq{{\textnormal{q}}}
\def\rr{{\textnormal{r}}}
\def\rs{{\textnormal{s}}}
\def\rt{{\textnormal{t}}}
\def\ru{{\textnormal{u}}}
\def\rv{{\textnormal{v}}}
\def\rw{{\textnormal{w}}}
\def\rx{{\textnormal{x}}}
\def\ry{{\textnormal{y}}}
\def\rz{{\textnormal{z}}}

% Random vectors
\def\rvepsilon{{\mathbf{\epsilon}}}
\def\rvtheta{{\mathbf{\theta}}}
\def\rva{{\mathbf{a}}}
\def\rvb{{\mathbf{b}}}
\def\rvc{{\mathbf{c}}}
\def\rvd{{\mathbf{d}}}
\def\rve{{\mathbf{e}}}
\def\rvf{{\mathbf{f}}}
\def\rvg{{\mathbf{g}}}
\def\rvh{{\mathbf{h}}}
\def\rvu{{\mathbf{i}}}
\def\rvj{{\mathbf{j}}}
\def\rvk{{\mathbf{k}}}
\def\rvl{{\mathbf{l}}}
\def\rvm{{\mathbf{m}}}
\def\rvn{{\mathbf{n}}}
\def\rvo{{\mathbf{o}}}
\def\rvp{{\mathbf{p}}}
\def\rvq{{\mathbf{q}}}
\def\rvr{{\mathbf{r}}}
\def\rvs{{\mathbf{s}}}
\def\rvt{{\mathbf{t}}}
\def\rvu{{\mathbf{u}}}
\def\rvv{{\mathbf{v}}}
\def\rvw{{\mathbf{w}}}
\def\rvx{{\mathbf{x}}}
\def\rvy{{\mathbf{y}}}
\def\rvz{{\mathbf{z}}}

% Elements of random vectors
\def\erva{{\textnormal{a}}}
\def\ervb{{\textnormal{b}}}
\def\ervc{{\textnormal{c}}}
\def\ervd{{\textnormal{d}}}
\def\erve{{\textnormal{e}}}
\def\ervf{{\textnormal{f}}}
\def\ervg{{\textnormal{g}}}
\def\ervh{{\textnormal{h}}}
\def\ervi{{\textnormal{i}}}
\def\ervj{{\textnormal{j}}}
\def\ervk{{\textnormal{k}}}
\def\ervl{{\textnormal{l}}}
\def\ervm{{\textnormal{m}}}
\def\ervn{{\textnormal{n}}}
\def\ervo{{\textnormal{o}}}
\def\ervp{{\textnormal{p}}}
\def\ervq{{\textnormal{q}}}
\def\ervr{{\textnormal{r}}}
\def\ervs{{\textnormal{s}}}
\def\ervt{{\textnormal{t}}}
\def\ervu{{\textnormal{u}}}
\def\ervv{{\textnormal{v}}}
\def\ervw{{\textnormal{w}}}
\def\ervx{{\textnormal{x}}}
\def\ervy{{\textnormal{y}}}
\def\ervz{{\textnormal{z}}}

% Random matrices
\def\rmA{{\mathbf{A}}}
\def\rmB{{\mathbf{B}}}
\def\rmC{{\mathbf{C}}}
\def\rmD{{\mathbf{D}}}
\def\rmE{{\mathbf{E}}}
\def\rmF{{\mathbf{F}}}
\def\rmG{{\mathbf{G}}}
\def\rmH{{\mathbf{H}}}
\def\rmI{{\mathbf{I}}}
\def\rmJ{{\mathbf{J}}}
\def\rmK{{\mathbf{K}}}
\def\rmL{{\mathbf{L}}}
\def\rmM{{\mathbf{M}}}
\def\rmN{{\mathbf{N}}}
\def\rmO{{\mathbf{O}}}
\def\rmP{{\mathbf{P}}}
\def\rmQ{{\mathbf{Q}}}
\def\rmR{{\mathbf{R}}}
\def\rmS{{\mathbf{S}}}
\def\rmT{{\mathbf{T}}}
\def\rmU{{\mathbf{U}}}
\def\rmV{{\mathbf{V}}}
\def\rmW{{\mathbf{W}}}
\def\rmX{{\mathbf{X}}}
\def\rmY{{\mathbf{Y}}}
\def\rmZ{{\mathbf{Z}}}

% Elements of random matrices
\def\ermA{{\textnormal{A}}}
\def\ermB{{\textnormal{B}}}
\def\ermC{{\textnormal{C}}}
\def\ermD{{\textnormal{D}}}
\def\ermE{{\textnormal{E}}}
\def\ermF{{\textnormal{F}}}
\def\ermG{{\textnormal{G}}}
\def\ermH{{\textnormal{H}}}
\def\ermI{{\textnormal{I}}}
\def\ermJ{{\textnormal{J}}}
\def\ermK{{\textnormal{K}}}
\def\ermL{{\textnormal{L}}}
\def\ermM{{\textnormal{M}}}
\def\ermN{{\textnormal{N}}}
\def\ermO{{\textnormal{O}}}
\def\ermP{{\textnormal{P}}}
\def\ermQ{{\textnormal{Q}}}
\def\ermR{{\textnormal{R}}}
\def\ermS{{\textnormal{S}}}
\def\ermT{{\textnormal{T}}}
\def\ermU{{\textnormal{U}}}
\def\ermV{{\textnormal{V}}}
\def\ermW{{\textnormal{W}}}
\def\ermX{{\textnormal{X}}}
\def\ermY{{\textnormal{Y}}}
\def\ermZ{{\textnormal{Z}}}

% Vectors
\def\vzero{{\bm{0}}}
\def\vone{{\bm{1}}}
\def\vmu{{\bm{\mu}}}
\def\vtheta{{\bm{\theta}}}
\def\va{{\bm{a}}}
\def\vb{{\bm{b}}}
\def\vc{{\bm{c}}}
\def\vd{{\bm{d}}}
\def\ve{{\bm{e}}}
\def\vf{{\bm{f}}}
\def\vg{{\bm{g}}}
\def\vh{{\bm{h}}}
\def\vi{{\bm{i}}}
\def\vj{{\bm{j}}}
\def\vk{{\bm{k}}}
\def\vl{{\bm{l}}}
\def\vm{{\bm{m}}}
\def\vn{{\bm{n}}}
\def\vo{{\bm{o}}}
\def\vp{{\bm{p}}}
\def\vq{{\bm{q}}}
\def\vr{{\bm{r}}}
\def\vs{{\bm{s}}}
\def\vt{{\bm{t}}}
\def\vu{{\bm{u}}}
\def\vv{{\bm{v}}}
\def\vw{{\bm{w}}}
\def\vx{{\bm{x}}}
\def\vy{{\bm{y}}}
\def\vz{{\bm{z}}}

% Elements of vectors
\def\evalpha{{\alpha}}
\def\evbeta{{\beta}}
\def\evepsilon{{\epsilon}}
\def\evlambda{{\lambda}}
\def\evomega{{\omega}}
\def\evmu{{\mu}}
\def\evpsi{{\psi}}
\def\evsigma{{\sigma}}
\def\evtheta{{\theta}}
\def\eva{{a}}
\def\evb{{b}}
\def\evc{{c}}
\def\evd{{d}}
\def\eve{{e}}
\def\evf{{f}}
\def\evg{{g}}
\def\evh{{h}}
\def\evi{{i}}
\def\evj{{j}}
\def\evk{{k}}
\def\evl{{l}}
\def\evm{{m}}
\def\evn{{n}}
\def\evo{{o}}
\def\evp{{p}}
\def\evq{{q}}
\def\evr{{r}}
\def\evs{{s}}
\def\evt{{t}}
\def\evu{{u}}
\def\evv{{v}}
\def\evw{{w}}
\def\evx{{x}}
\def\evy{{y}}
\def\evz{{z}}

% Matrix
\def\mA{{\bm{A}}}
\def\mB{{\bm{B}}}
\def\mC{{\bm{C}}}
\def\mD{{\bm{D}}}
\def\mE{{\bm{E}}}
\def\mF{{\bm{F}}}
\def\mG{{\bm{G}}}
\def\mH{{\bm{H}}}
\def\mI{{\bm{I}}}
\def\mJ{{\bm{J}}}
\def\mK{{\bm{K}}}
\def\mL{{\bm{L}}}
\def\mM{{\bm{M}}}
\def\mN{{\bm{N}}}
\def\mO{{\bm{O}}}
\def\mP{{\bm{P}}}
\def\mQ{{\bm{Q}}}
\def\mR{{\bm{R}}}
\def\mS{{\bm{S}}}
\def\mT{{\bm{T}}}
\def\mU{{\bm{U}}}
\def\mV{{\bm{V}}}
\def\mW{{\bm{W}}}
\def\mX{{\bm{X}}}
\def\mY{{\bm{Y}}}
\def\mZ{{\bm{Z}}}
\def\mBeta{{\bm{\beta}}}
\def\mPhi{{\bm{\Phi}}}
\def\mLambda{{\bm{\Lambda}}}
\def\mSigma{{\bm{\Sigma}}}

% Tensor
\DeclareMathAlphabet{\mathsfit}{\encodingdefault}{\sfdefault}{m}{sl}
\SetMathAlphabet{\mathsfit}{bold}{\encodingdefault}{\sfdefault}{bx}{n}
\newcommand{\tens}[1]{\bm{\mathsfit{#1}}}
\def\tA{{\tens{A}}}
\def\tB{{\tens{B}}}
\def\tC{{\tens{C}}}
\def\tD{{\tens{D}}}
\def\tE{{\tens{E}}}
\def\tF{{\tens{F}}}
\def\tG{{\tens{G}}}
\def\tH{{\tens{H}}}
\def\tI{{\tens{I}}}
\def\tJ{{\tens{J}}}
\def\tK{{\tens{K}}}
\def\tL{{\tens{L}}}
\def\tM{{\tens{M}}}
\def\tN{{\tens{N}}}
\def\tO{{\tens{O}}}
\def\tP{{\tens{P}}}
\def\tQ{{\tens{Q}}}
\def\tR{{\tens{R}}}
\def\tS{{\tens{S}}}
\def\tT{{\tens{T}}}
\def\tU{{\tens{U}}}
\def\tV{{\tens{V}}}
\def\tW{{\tens{W}}}
\def\tX{{\tens{X}}}
\def\tY{{\tens{Y}}}
\def\tZ{{\tens{Z}}}


% Graph
\def\gA{{\mathcal{A}}}
\def\gB{{\mathcal{B}}}
\def\gC{{\mathcal{C}}}
\def\gD{{\mathcal{D}}}
\def\gE{{\mathcal{E}}}
\def\gF{{\mathcal{F}}}
\def\gG{{\mathcal{G}}}
\def\gH{{\mathcal{H}}}
\def\gI{{\mathcal{I}}}
\def\gJ{{\mathcal{J}}}
\def\gK{{\mathcal{K}}}
\def\gL{{\mathcal{L}}}
\def\gM{{\mathcal{M}}}
\def\gN{{\mathcal{N}}}
\def\gO{{\mathcal{O}}}
\def\gP{{\mathcal{P}}}
\def\gQ{{\mathcal{Q}}}
\def\gR{{\mathcal{R}}}
\def\gS{{\mathcal{S}}}
\def\gT{{\mathcal{T}}}
\def\gU{{\mathcal{U}}}
\def\gV{{\mathcal{V}}}
\def\gW{{\mathcal{W}}}
\def\gX{{\mathcal{X}}}
\def\gY{{\mathcal{Y}}}
\def\gZ{{\mathcal{Z}}}

% Sets
\def\sA{{\mathbb{A}}}
\def\sB{{\mathbb{B}}}
\def\sC{{\mathbb{C}}}
\def\sD{{\mathbb{D}}}
% Don't use a set called E, because this would be the same as our symbol
% for expectation.
\def\sF{{\mathbb{F}}}
\def\sG{{\mathbb{G}}}
\def\sH{{\mathbb{H}}}
\def\sI{{\mathbb{I}}}
\def\sJ{{\mathbb{J}}}
\def\sK{{\mathbb{K}}}
\def\sL{{\mathbb{L}}}
\def\sM{{\mathbb{M}}}
\def\sN{{\mathbb{N}}}
\def\sO{{\mathbb{O}}}
\def\sP{{\mathbb{P}}}
\def\sQ{{\mathbb{Q}}}
\def\sR{{\mathbb{R}}}
\def\sS{{\mathbb{S}}}
\def\sT{{\mathbb{T}}}
\def\sU{{\mathbb{U}}}
\def\sV{{\mathbb{V}}}
\def\sW{{\mathbb{W}}}
\def\sX{{\mathbb{X}}}
\def\sY{{\mathbb{Y}}}
\def\sZ{{\mathbb{Z}}}

% Entries of a matrix
\def\emLambda{{\Lambda}}
\def\emA{{A}}
\def\emB{{B}}
\def\emC{{C}}
\def\emD{{D}}
\def\emE{{E}}
\def\emF{{F}}
\def\emG{{G}}
\def\emH{{H}}
\def\emI{{I}}
\def\emJ{{J}}
\def\emK{{K}}
\def\emL{{L}}
\def\emM{{M}}
\def\emN{{N}}
\def\emO{{O}}
\def\emP{{P}}
\def\emQ{{Q}}
\def\emR{{R}}
\def\emS{{S}}
\def\emT{{T}}
\def\emU{{U}}
\def\emV{{V}}
\def\emW{{W}}
\def\emX{{X}}
\def\emY{{Y}}
\def\emZ{{Z}}
\def\emSigma{{\Sigma}}

% entries of a tensor
% Same font as tensor, without \bm wrapper
\newcommand{\etens}[1]{\mathsfit{#1}}
\def\etLambda{{\etens{\Lambda}}}
\def\etA{{\etens{A}}}
\def\etB{{\etens{B}}}
\def\etC{{\etens{C}}}
\def\etD{{\etens{D}}}
\def\etE{{\etens{E}}}
\def\etF{{\etens{F}}}
\def\etG{{\etens{G}}}
\def\etH{{\etens{H}}}
\def\etI{{\etens{I}}}
\def\etJ{{\etens{J}}}
\def\etK{{\etens{K}}}
\def\etL{{\etens{L}}}
\def\etM{{\etens{M}}}
\def\etN{{\etens{N}}}
\def\etO{{\etens{O}}}
\def\etP{{\etens{P}}}
\def\etQ{{\etens{Q}}}
\def\etR{{\etens{R}}}
\def\etS{{\etens{S}}}
\def\etT{{\etens{T}}}
\def\etU{{\etens{U}}}
\def\etV{{\etens{V}}}
\def\etW{{\etens{W}}}
\def\etX{{\etens{X}}}
\def\etY{{\etens{Y}}}
\def\etZ{{\etens{Z}}}

% The true underlying data generating distribution
\newcommand{\pdata}{p_{\rm{data}}}
% The empirical distribution defined by the training set
\newcommand{\ptrain}{\hat{p}_{\rm{data}}}
\newcommand{\Ptrain}{\hat{P}_{\rm{data}}}
% The model distribution
\newcommand{\pmodel}{p_{\rm{model}}}
\newcommand{\Pmodel}{P_{\rm{model}}}
\newcommand{\ptildemodel}{\tilde{p}_{\rm{model}}}
% Stochastic autoencoder distributions
\newcommand{\pencode}{p_{\rm{encoder}}}
\newcommand{\pdecode}{p_{\rm{decoder}}}
\newcommand{\precons}{p_{\rm{reconstruct}}}

\newcommand{\laplace}{\mathrm{Laplace}} % Laplace distribution

\newcommand{\E}{\mathbb{E}}
\newcommand{\Ls}{\mathcal{L}}
\newcommand{\R}{\mathbb{R}}
\newcommand{\emp}{\tilde{p}}
\newcommand{\lr}{\alpha}
\newcommand{\reg}{\lambda}
\newcommand{\rect}{\mathrm{rectifier}}
\newcommand{\softmax}{\mathrm{softmax}}
\newcommand{\sigmoid}{\sigma}
\newcommand{\softplus}{\zeta}
\newcommand{\KL}{D_{\mathrm{KL}}}
\newcommand{\Var}{\mathrm{Var}}
\newcommand{\standarderror}{\mathrm{SE}}
\newcommand{\Cov}{\mathrm{Cov}}
% Wolfram Mathworld says $L^2$ is for function spaces and $\ell^2$ is for vectors
% But then they seem to use $L^2$ for vectors throughout the site, and so does
% wikipedia.
\newcommand{\normlzero}{L^0}
\newcommand{\normlone}{L^1}
\newcommand{\normltwo}{L^2}
\newcommand{\normlp}{L^p}
\newcommand{\normmax}{L^\infty}

\newcommand{\parents}{Pa} % See usage in notation.tex. Chosen to match Daphne's book.

\DeclareMathOperator*{\argmax}{arg\,max}
\DeclareMathOperator*{\argmin}{arg\,min}

\DeclareMathOperator{\sign}{sign}
\DeclareMathOperator{\Tr}{Tr}
\let\ab\allowbreak

   % Pour l'environnement proof

\usepackage{hyperref}
\hypersetup{colorlinks,linkcolor={blue},citecolor={blue},urlcolor={red}}
\usepackage{url}
\usepackage{algorithm}
\usepackage{algpseudocode}

\usepackage{float} % Pour utiliser l'option H pour les tables
\usepackage{caption}
\captionsetup[table]{position=above}

\usepackage{amsmath}
\usepackage{amssymb}
\usepackage{amsthm}
\usepackage{enumitem}

\usepackage{todonotes}
\usepackage{multirow}
\newtheorem{theorem}{Theorem}
\newtheorem{lemma}{Lemma}
\newtheorem{Remark}{Remark}
\newtheorem{assumption}{Assumption}
\usepackage{cleveref}
\usepackage{graphicx}
\crefname{assumption}{Assumption}{Assumptions}
\crefname{Remark}{Remark}{Remark}

\title{On the Almost Sure Convergence of the Stochastic Three Points Algorithm}

% Authors must not appear in the submitted version. They should be hidden
% as long as the \iclrfinalcopy macro remains commented out below.
% Non-anonymous submissions will be rejected without review.

\author{Taha EL BAKKALI EL KADI, Omar SAADI \\
UM6P College of Computing\\
Benguerir, Morocco \\
\texttt{\{taha.elbakkali,omar.saadi\}@um6p.ma} 
}

% Do not uncomment the following line for preprint submission
%\iclrfinalcopy % Uncomment for camera-ready version, but NOT for submission.
 \iclrfinalcopy
\begin{document}

\maketitle


\begin{abstract}
 The stochastic three points (STP) algorithm is a derivative-free optimization technique designed for unconstrained optimization problems in $\mathbb{R}^d$. In this paper, we analyze this algorithm for three classes of functions: smooth functions that may lack convexity, smooth convex functions, and smooth functions that are strongly convex. Our work provides the first almost sure convergence results of the STP algorithm, alongside some convergence results in expectation.
For the class of smooth functions, we establish that the best gradient iterate of the STP algorithm converges almost surely to zero at a rate of $o(1/{T^{\frac{1}{2}-\epsilon}})$ for any $\epsilon\in (0,\frac{1}{2})$, where $T$ is the number of iterations. Furthermore, within the same class of functions, we establish both almost sure convergence and convergence in expectation of the final gradient iterate towards zero.
For the class of smooth convex functions, we establish that $f(\theta^T)$ converges to $\inf_{\theta \in \mathbb{R}^d} f(\theta)$ almost surely at a rate of $o(1/{T^{1-\epsilon}})$ for any $\epsilon\in (0,1)$, and in expectation at a rate of $O(\frac{d}{T})$ where $d$ is the dimension of the space.
Finally, for the class of smooth functions that are strongly convex, we establish that when step sizes are obtained by approximating the directional derivatives of the function, $f(\theta^T)$ converges to $\inf_{\theta \in \mathbb{R}^d} f(\theta)$ in expectation at a rate of $O((1-\frac{\mu}{2\pi dL})^T)$, and almost surely at a rate of $o((1-s\frac{\mu}{2\pi dL})^T)$ for any $s\in (0,1)$,  where $\mu$ and $L$
are  the strong convexity and smoothness parameters of the function.
\end{abstract}


\section{Introduction}
We are interested in the minimization of a smooth function $f : \mathbb{R}^d \mapsto \mathbb{R}$:
%following optimization problem :
$$\min_{\theta \in \mathbb{R}^d} f(\theta),$$
where we work within the constraint of not having access to the derivatives of 
$f$, relying exclusively on a function evaluation oracle.
%the objective function $f : \mathbb{R}^d \mapsto \mathbb{R} $ is differentiable and bounded from below. In this context,
 The methods used %to solve this minimization problem
 in this framework are called derivative-free methods or zeroth-order methods \citep{conn2009introduction,ghadimi2013stochastic,nesterov2017random,larson2019derivative,golovin2019gradientless,bergou2020stochastic}. They are increasingly embraced for solving many machine learning problems where obtaining gradient information is either impractical or computationally expensive,  %These methods 
 remaining crucial in applications such as generating adversarial attacks on deep neural network classifiers \citep{chen2017zoo,tu2019autozoom}, reinforcement learning \citep{malik2019derivative,salimans2017evolution}, and hyperparameter tuning of ML models \citep{snoek2012practical,turner2021bayesian}. Therefore, exploring the theoretical properties of derivative-free methods is not only of theoretical interest but also crucial for practical applications.


Zeroth-order optimization methods can be divided into two main categories: direct search methods and gradient estimation methods. In direct search methods, the objective function is evaluated along a set of directions to guarantee descent by taking appropriate small step sizes. These directions can be either deterministic \citep{vicente2013worst} or stochastic \citep{golovin2019gradientless,bergou2020stochastic}. In contrast, gradient estimation methods approximate the gradient of the objective function using zeroth-order information to design approximate gradient methods \citep{nesterov2017random,shamir2017optimal}.

A recent and noteworthy zeroth-order method is the Stochastic Three Points (STP) algorithm (see \Cref{Alg1}), a directed search method with stochastic search directions, introduced by \cite{bergou2020stochastic}. The STP algorithm stands out among zeroth-order methods for its balance of simplicity and strong theoretical guarantees.
\begin{algorithm}[H]
  \caption{Stochastic Three Points (STP)}
  \label{alg:stp}\label{Alg1}
  \begin{algorithmic}[1]
    \State \textbf{Input:} $\theta^1 \in \mathbb{R}^d$, step size sequence $\{\alpha_t\}_{t \ge 1} \in  (0,\infty)^{\mathbb{N}^{*} }$, probability distribution $\mathcal{D}$ on $\mathbb{R}^d$.
    \For{$t=1,2, \ldots$}
      \State Generate a random vector $s_t \sim \mathcal{D}$, 
      \State $\theta^{t+1} = \argmin_{\theta  \in \{  \theta^t, \theta^t + \alpha_t s_t,\theta^t - \alpha_t s_t \} } f(\theta)$.
    \EndFor
  \end{algorithmic}
\end{algorithm}


Compared to deterministic directed search (DDS) methods, the worst-case complexity bounds for STP are similar; however, they differ in their dependence on the problem's dimensionality. For STP, the bounds increase linearly with the dimension \citep{bergou2020stochastic}, whereas for DDS, they increase quadratically \citep{konevcny2014simple,vicente2013worst}. Specifically, when the objective function is smooth, STP requires \( O(d \epsilon^{-2}) \) function evaluations to get a gradient with norm smaller than \( \epsilon \), in expectation. For smooth, convex functions with a minimum and a bounded sublevel set, the complexity is \( O(d \epsilon^{-1}) \) to find an \( \epsilon\)-optimal solution. In the strongly convex case, this complexity reduces further to \( O(d \log \epsilon^{-1}) \). In all these cases, DDS methods exhibit analogous complexity bounds but with a quadratic dependence on \( d \), i.e., \( d^2 \) instead of \( d \). 
In comparison to directed search with stochastic directions, STP also matches the complexity bound derived by Gratton et al. \citep{doi:10.1137/140961602} for the smooth case, which is the only case they address in their work. In their approach, a decrease condition is imposed to determine whether to accept or reject a step based on a set of random directions. The Gradientless Descent (GLD) algorithm \citep{golovin2019gradientless} is another direct search method with stochastic directions. Golovin et al. show that an $\epsilon$-optimal solution can be found in $O(kQ \log(d) \log\left(R \epsilon^{-1}\right))$ for any monotone transform of a smooth and strongly convex function with latent dimension $k < d$, where the input dimension is $d$, $R$ is the diameter of the input space, and $Q$ is the condition number. When the monotone transformation is the identity and $k = d$, this complexity is higher than the one obtained for the STP algorithm by a factor of $\log(d)$. However, it is important to note that monotone transforms of smooth and strongly convex functions are not necessarily strongly convex.

Compared to approximate gradient methods, STP matches the complexity bounds of the random gradient-free (RGF) 
algorithm \citep{nesterov2017random} (see \cref{exper}) across the three cases: smooth non-convex, smooth convex, and smooth strongly convex. 
This matching in complexities is in terms of the accuracy $\epsilon$ and the dimensionality $d$.  %Note that the rate obtained for STP in the convex case needs the additional assumption that there  the sublevel set of the initial point is bounded. In fact, in this convex case when the search directions are sampled among the canonical vectors of $\mathbb{R}^d$, STP can be seen as a zeroth order version of the randomized coordinate gradient descent algorithm, where the same assumption is considered, see~\cite[Assumption 1]{wright2015coordinate}.

%\todo{Mention Optimal ZO rates (Accelerated versions of RGF and STP)}




In practical terms, for classical applications of zeroth-order methods, STP variants demonstrate strong performance when compared to state-of-the-art methods. For instance, in reinforcement learning and continuous control, specifically in the MuJoCo simulation suite~\citep{6386109}, STP with momentum (which, in expectation, achieves the same complexity bounds as standard STP, see~\cite{gorbunov2020momentum}) outperforms methods like Augmented Random Search (ARS), Trust Region Policy Optimization (TRPO), and Natural Policy Gradient (NG) across environments such as Swimmer-v1, Hopper-v1, HalfCheetah-v1, and Ant-v1. Even in the more challenging Humanoid-v1 environment, STP with momentum achieves competitive results~\citep{gorbunov2020momentum}. Additionally, in the context of generating adversarial attacks on deep neural network classifiers, the Minibatch Stochastic Three Points (MiSTP)  method ~\citep{boucherouite2024minibatch} demonstrates superior performance compared to other variants of zero-order methods, that are adapted to the stochastic setting, such as RSGF (also called ZO-SGD) \citep{ghadimi2013stochastic}, ZO-SVRG-Ave, and ZO-SVRG \citep{liu2018zeroth}.


Within the realm of first-order optimization methods that rely on gradient information, numerous studies have investigated the almost sure convergence of the Stochastic Gradient Descent (SGD)  algorithm and its variants \citep{bertsekas2000gradient, nguyen2019new, mertikopoulos2020almost, sebbouh2021almost, liu2022almost}. In contrast, the literature on the almost sure convergence of zeroth-order methods remains less developed compared to that of SGD. 

In \citep{doi:10.1137/140961602}, the authors investigate zeroth-order direct-search methods under a probabilistic descent framework. Specifically, they generate randomly the search directions, while assuming that with a certain probability at least one of them is of descent type. For smooth objective functions, their analysis establishes (in Theorem $3.4$) the almost sure convergence of the best iterate of the gradient norm to zero. However, the analysis does not provide a convergence rate for this almost sure convergence result, nor does it guarantee the convergence of the gradient norm of the last iterate. In our paper, we provide such results for the STP algorithm (see~\Cref{Table1}). 
\citep{doi:10.1137/140961602} also establish (in Corollary $4.7$) a convergence rate $O(1/\sqrt{T})$ for the best iterate with overwhelmingly high probability, but this rate is still not guaranteed almost surely. In our work, we provide the first almost sure convergence rate of the best iterate for zeroth-order methods (see~\Cref{Table1}). 
More recently, \cite{wang2022convergence} explore the convergence of the Stochastic Zeroth-order Gradient Descent (SZGD) algorithm for objective functions satisfying the \L{}ojasiewicz inequality. Assuming smoothness, they demonstrated (in Lemma $1$) that the gradient norm of the last iterate converges to zero. Furthermore, in Lemma $2$, they proved that the sequence generated by the SZGD algorithm converges almost surely to a critical point, which is a stronger result, since the gradient of $f$ is continuous. However, this analysis is limited to \L{}ojasiewicz functions, which, by definition, satisfy a strong property which is the property essentially used in the analysis of strongly convex functions.

In this paper, we are interested in studying the almost sure convergence of the STP algorithm. For the three classes of functions (smooth, smooth convex, and smooth strongly convex),
first convergence results, in terms of expectation, were provided in \cite{bergou2020stochastic}. However,
it is crucial to note that ensuring almost sure convergence properties is essential for understanding
the behavior of each trajectory of the STP algorithm and guaranteeing that any instantiation of the
algorithm converges with probability one.




%In X, they provide an almost sure convergence of the last iterate of a trust region based zero-order method. This method is based on iteratively approximating the objective function using quadratic functions.   




 
%However, the theoretical analysis of this algorithm is conducted under a specific setting. 

%We propose to compare both algorithms in the simulation section, specifically in the context of non-convex optimization.






%A deterministic variant of STP is the Compass Search (CS) algorithm \citep{audet2014survey}. 
%It starts with an initial solution \(\theta^1\) and a step size parameter \(\Delta_1\). 
%At each iteration, the algorithm generates trial points by adjusting the coordinates of the current vector \(\theta^t\) by \(\pm \Delta_t\). 
%If a trial point improves the objective function, it becomes the updated vector $\theta^{t+1}$, with the step size remaining unchanged. 
%Otherwise, the current vector is retained, and the step size \(\Delta_t\) is halved. 
%In contrast, the CS algorithm may require \(O(d)\) function evaluations per iteration in the worst case when assessing all coordinate directions \( \pm e_i\) (\(e_i\) are the canonical basis vectors in \(\mathbb{R}^d\)), while the STP algorithm performs only two new function evaluations at each iteration.







\textbf{Our Contribution $\&$ Related Work.}
In cases where the only verified assumptions regarding the function are its smoothness and having a lower bound, Bergou et al. established in their paper \cite[Theorem 4.1]{bergou2020stochastic} that by using \Cref{Alg1} and selecting a step size sequence $\{\frac{\alpha}{\sqrt{t}}\}_{t \ge 1 }$ with $\alpha>0$, the best gradient iterate converges in expectation to $0$ at a rate of $O(\frac{\sqrt{d}}{\sqrt{T}}).$  Expanding on this, we prove that employing a similar step size sequence $\{\frac{\alpha}{t^{\frac{1}{2}+\epsilon} }\}_{t \ge 1}$ with $\epsilon \in (0,\frac{1}{2})$  results in an almost sure convergence rate of $o(\frac{1}{T^{\frac{1}{2}-\epsilon }} )  $, which is arbitrarily close to the rate achieved for the convergence in expectation when $\epsilon$ is close to $0$ (see ~\Cref{Thm1}). It's worth noting that a similar almost sure convergence result has been established for the SGD Algorithm. For more information, refer to \cite[Corollary 18]{sebbouh2021almost} and   \cite[Theorem 1]{liu2022almost}. However, it should be noted that this similar result for the SGD Algorithm is provided for $\min_{1 \le t \le T} \lVert \nabla f(\theta^t) \rVert^2$, while for the STP Algorithm, it is provided for $\min_{1 \le t \le T} \lVert \nabla f(\theta^t) \rVert$.  More precisely, for the STP algorithm, we have
$\min_{1 \le t \le T} \lVert \nabla f(\theta^t) \rVert=o(1/T^{ \frac{1}{2} - \epsilon})$, while for the SGD algorithm, we have $\min_{1 \le t \le T} \lVert \nabla f(\theta^t) \rVert=o(1/T^{ \frac{1}{4} - \frac{\epsilon}{2}} )$. The issue with both convergence results, whether it's the one by Bergou et al. \cite[Theorem 4.1]{bergou2020stochastic}  about the convergence in expectation or our first result about the almost sure convergence, is that they don't guarantee the gradient of $f$ at the final point $\theta^T$ to be small (either in expectation or almost surely). Instead, they assure that the gradient of $f$ at some point produced by the STP algorithm is small. In our paper, we additionally prove that the gradient of $f$ at the final point $\theta^T$ converges to $0$ almost surely and in expectation without requiring additional assumptions about the function beyond its smoothness and having a lower bound (see~\Cref{Thm2,Thm3}). Notably, for the case of the SGD algorithm, the question of the almost sure convergence of the last gradient iterate has been addressed in \cite{bertsekas2000gradient}.

For smooth convex functions, if $f$ has a global minimum $\theta^{*}$ and possesses a bounded sublevel set, we show that selecting a step size sequence $\alpha_t=O(\frac{1}{t^{1-\beta}})$ for some $\beta \in (0 ,\frac{1}{2})$ ensures that $f(\theta^T)$ converges almost surely to $ f(\theta^{*})$ at a rate of $o(\frac{1}{T^{1-\epsilon}})$ for all $\epsilon \in (2\beta,1)$ (see ~\Cref{Thm5}). A similar result, with the same convergence rate and the same criteria for choosing the step size sequence, is established for the stochastic Nesterov's accelerated gradient algorithm by Jun Liu et al. in \cite[Theorem 3]{liu2022almost}. For the same class of functions and under the same assumptions, Bergou et al. established in \cite[Theorem 5.5]{bergou2020stochastic} that for a fixed precision $\epsilon$ and a sufficiently large number of iterations $T$ on the order of $\frac{1}{\epsilon}$, by selecting a step size sequence $\{\frac{|f(\theta^t+hs_t )-f(\theta^t)|}{Lh} \}_{t \ge 1}$ where $h$ is sufficiently small on the order of $\frac{}{}\mathbb{E}[f(\theta^{T-1})]- f(\theta^{*})$,  one can get: $\mathbb{E}[f(\theta^{T})]-f(\theta^{*}) \le \epsilon$. 
Here, the choice of $h$ depends on the quantity $\mathbb{E}[f(\theta^{T-1})]$ which is not known at the begining.
%In this theorem, to define the step size sequence, one must choose $h$ to be sufficiently small, on the order of $\mathbb{E}[f(\theta^{T-1})]-\inf_{\theta \in \mathbb{R}^d } f(\theta)$. 
%However, defining $\theta^{T-1}$ itself depends on the choice of the step size sequence, which in turn depends on $h$. 
%This circular dependence makes it difficult to verify the condition on $h$ in practice. 
Moreover, the theorem does not guarantee that $\mathbb{E}[f(\theta^T)]$ converges to $f(\theta^{*})$, because the step sizes depend on $\epsilon$. In contrast, in \Cref{Thm4}, we show that by selecting a step size sequence $\{\frac{\alpha}{t}\}_{t \ge 1}$, where $\alpha$ is suitably chosen, $\mathbb{E}[f(\theta^T)]$ converges to $f(\theta^{*})$ at a rate of $O(\frac{d}{T})$  . 

For smooth, strongly convex functions, Bergou et al. established in \cite[Theorem 6.3]{bergou2020stochastic} that, for any $\epsilon > 0$, using the step size sequence $  \{\frac{|f(\theta^{t} + h s_t) - f(\theta^t)|}{Lh}\}_{t\ge 1}$, where $h$ is small on the order of $\sqrt{\epsilon}$, the gap between the expected value of the objective function and its infimum stays within $\epsilon$ accuracy for 
%a large enough $K$ 
a number of iterations 
on the order of $\log(\frac{ f(\theta^t) - \inf_{\theta \in \mathbb{R}^d } f(\theta) }{\epsilon})$. 
However, this result doesn't indicate how the gap $\mathbb{E}[f(\theta^T)] -\inf_{\theta \in \mathbb{R}^d} f(\theta) $ improves with more iterations and does not guarantee convergence since the step sizes depend on $\epsilon$. To address this issue, we define the step size sequence as $\{ \frac{|f(\theta^{t} + h^{-t} s_t) - f(\theta^t)|}{Lh^{-t}} \}_{t \ge 1}$ with a suitable $h$, leading to a convergence rate of $O((1-\frac{\mu}{2\pi dL})^T)$ in expectation, and $o((1-s \frac{\mu}{2\pi dL})^T)$ almost surely for all $s \in (0,1)$ (see ~\Cref{Thm6,Thm7}). All of our convergence rates are succinctly presented in~\Cref{Table1}.
%However, a limitation of this result arises due to machine precision when dividing by $h^{-t}$  to compute the step sizes. This precision limitation can impact the algorithm's convergence speed.  It is noted that if $f$ is strongly convex, then for all vector $c \in \mathbb{R}^d$, the sublevel set of $f$ at $c$, $L(c) = \{\theta \in \mathbb{R}^d \mid f(\theta) \le f(c) \}$ is bounded (see \Cref{Rem8}), and there exists $\theta^{*}$ such that $f(\theta^{*})= \inf_{\theta \in \mathbb{R}^d} f(\theta)$.Therefore, we can obtain the same results as \Cref{Thm4,Thm5} in the case where the function is smooth strongly convex, ie : by selecting a step size sequence $\{\frac{\alpha}{t}\}_{t \ge 1}$, where $\alpha$ is suitably chosen to be sufficiently large, we can ensure that the gradient iterate of the STP algorithm converges in expectation at a rate of $O(\frac{1}{T})$.
\begin{table}[h!]
\caption{Summary of convergence rates for the STP algorithm.}
\centering
\fontsize{5.83pt}{9.2pt}\selectfont % Font size and baselineskip
\begin{tabular}{cccccc}
\hline
\rowcolor[gray]{0.9}  % Gray color for the first row from the second column
\textbf{Functions} & \textbf{Assump} & \textbf{Step size} &  \textbf{Iterate} & \textbf{Conv / Rate} & \textbf{Ref} \\
\hline
\multirow{3}{*}{Smooth} & \ref{Ass1},\ref{Ass5},\ref{Ass6} &
\multirow{2}{*}{\shortstack{$\{\frac{\alpha}{\sqrt{t}} \}_{t \ge 1}, \alpha>0$\\ \text{ }}} &
$\underset{1 \leq t \leq T}{\min} \lVert \nabla f(\theta^t) \rVert$  
& $\mathbb{E}/O\left(\frac{\sqrt{d}}{\sqrt{T}}\right)$ & 
$\begin{array}{l}
\text{Thm 4.1}\\
\text{\cite{bergou2020stochastic}}
\end{array}$



\\

& \ref{Ass1},\ref{Ass5},\ref{Ass6} & 
$\begin{cases}
\{\frac{\alpha}{t^{\frac{1}{2}+\epsilon }}\}_{t \ge 1}, \alpha>0\\
\epsilon \in (0, \frac{1}{2})
\end{cases}$ &
$\underset{1 \leq t \leq T}{\min} \lVert \nabla f(\theta^t) \rVert$ &$\text{a.s.} \; / o(\frac{1}{T^{\frac{1}{2}- \epsilon }})$ & Thm \ref{Thm1} \\

& \ref{Ass1},\ref{Ass5},\ref{Ass6} & 
$\begin{cases}
\{\alpha_t\}_{t \ge 1}\\
\sum_{t=1}^{\infty} \alpha_t^2 <\infty\\
\sum_{t=1}^{\infty} \alpha_t=\infty 
\end{cases}$ &
$ \lVert \nabla f(\theta^T) \rVert$ &$\mathbb{E} \; \&$  \text{a.s.}  / $o\left(1\right)$ & 
$\begin{array}{l}
\text{Thm \ref{Thm2}}\\
\text{Thm \ref{Thm3}}
\end{array}$

\\
\hline

\multirow{2}{*}{\shortstack{Smooth,\\ convex}} & \ref{Ass1},\ref{Ass2},\ref{Ass3},\ref{Ass5},\ref{Ass6} &
$\alpha_t=\frac{\alpha}{t}, \alpha \text{ is suitably chosen}$ &
$ f(\theta^T)-f(\theta^*)$ & $\mathbb{E}/O\left(\frac{d}{T}\right)$ & Thm \ref{Thm4} \\

& \ref{Ass1},\ref{Ass2},\ref{Ass3},\ref{Ass5},\ref{Ass6} & 
$\alpha_t= O\left(\frac{1}{t^{1-\beta}} \right), \beta \in (0,\frac{1}{2})$ &
$f(\theta^T)-f(\theta^*)$  &
$\begin{array}{l}
\text{a.s.}\; / o\left(\frac{1}{T^{1- \epsilon }}\right), \\
\forall \epsilon \in (2 \beta,1)
\end{array}$ & Thm \ref{Thm5} \\
\hline

\multirow{2}{*}{\shortstack{Smooth,\\strongly\\ convex}} & 
\ref{Ass1},\ref{Ass4},\ref{Ass5},\ref{Ass6},\ref{Ass7} &
$\begin{cases}
\{ \frac{|f(\theta^{t} + h^{-t} s_t) - f(\theta^t)|}{Lh^{-t}} \}_{t \ge 1}\\
h \text{ is large enough}
\end{cases}$ &
$ f(\theta^T)-f(\theta^*)$ &
$\begin{array}{l}
\mathbb{E}/O( (1- \beta)^T ) \\
 \beta\sim \frac{\mu}{dL}
\end{array}$

 & Thm \ref{Thm7} \\

& \ref{Ass1},\ref{Ass4},\ref{Ass5},\ref{Ass6},\ref{Ass7} & 
$\begin{cases}
\{ \frac{|f(\theta^{t} + h^{-t} s_t) - f(\theta^t)|}{Lh^{-t}} \}_{t \ge 1}\\
h \text{ is large enough}
\end{cases}$ &
$ f(\theta^T)-f(\theta^*)$ & 
$\begin{array}{l}
\text{a.s.}\; / o((1- \beta)^T) \\
 \beta\sim \frac{s \mu}{dL}; \forall s \in (0,1)
\end{array}$
& Thm \ref{Thm6} \\
\hline
\end{tabular}
\label{Table1}
\end{table}
\section{Problem setup and assumptions}\label{sec2}
We are interested in the following optimization problem:
$$\min_{\theta \in \mathbb{R}^d} f(\theta),$$
where the objective function $f : \mathbb{R}^d \mapsto \mathbb{R} $ is differentiable and bounded from below. In this context,
we work within the constraint of not having access to the derivatives of 
$f$, relying exclusively on a function evaluation oracle. 

Throughout the rest of the paper, we assume that the objective function is differentiable and bounded from below. We consider the following additional assumptions about $f$:
\begin{assumption}\label{Ass1}
$f$ is $L-$smooth, i.e., 
$\forall x,y \in \mathbb{R}^d,\; || \nabla f(x)-\nabla f(y) ||_2 \le  L \lVert x-y \rVert_2.$
\end{assumption}

Note that \cref{Ass1}, implies the following result \cite[Lemma 1.2.3]{nesterov2013introductory}: 
\begin{equation}\label{ineq-smooth}
\forall x,y \in \mathbb{R}^d,\; |f(y) - f(x)- \langle \nabla f(x),y-x \rangle | \le \frac{L}{2} \lVert y-x\rVert_{2}^2.
\end{equation}
\begin{assumption}\label{Ass2}
 $\exists \theta^{*}  \in \mathbb{R}^d, \; f(\theta^{*})=\inf_{\theta \in \mathbb{R}^d } f(\theta)$.
\end{assumption}


\begin{assumption}\label{Ass3}
%$f$ is convex, i.e. $(\forall x,y \in \mathbb{R}^d) : f(y) \ge f(x)+\langle \nabla f(x),y-x \rangle$
 $f$ is convex and there exists $ c \in \mathbb{R}^d$ such that the sublevel set of $f$ defined by $c$ is bounded, i.e.,
\begin{enumerate}
    \item $\forall x,y \in \mathbb{R}^d,\;  f(y) \ge f(x)+\langle \nabla f(x),y-x \rangle$.
    \item There exists $c \in \mathbb{R}^d$ such that $L( c)=\{x \in \mathbb{R}^d  \mid f(x) \leq f( c) \}$ is bounded.
\end{enumerate} 
\end{assumption}
\begin{assumption}\label{Ass4}
$f$ is $\mu$-strongly convex, i.e., 
there exists a positive constant $\mu$ such that: 
$$\forall x,y \in \mathbb{R}^d,\; f(y) \ge f(x)+ \langle \nabla f(x),y-x \rangle + \frac{\mu}{2} \lVert y-x\rVert_2^2.$$    
\end{assumption}

Note that  \cref{Ass4}, implies the following result \cite[Theorem 2.1.8]{nesterov2013introductory}: 
$$\forall x \in \mathbb{R}^d,\; \frac{1}{2 \mu} \lVert \nabla f(x) \rVert_2^2 \ge f(x)-\inf_{y \in \mathbb{R}^d} f(y)\;\;\;\; \text{(Polyak-{\L}ojasiewicz inequality)}.$$
For the distributions $\mathcal{D}$ over $\mathbb{R}^d$, we make the following assumptions:  
\begin{assumption}\label{Ass5}
    The probability distribution $\mathcal{D}$ on $\mathbb{R}^d$ satisfies:
\begin{enumerate}
    \item $\gamma_{\mathcal{D}}:=\mathbb{E}_{s \sim \mathcal{D}}[ \lVert s \rVert_2^2 ]<\infty$.
    \item There exists a norm $\lVert .  \rVert_{\mathcal{D}}$ on $\mathbb{R}^d$, for which we can find a  constant $\mu_\mathcal{D}>0$ such that:
    %There exists a constant $\mu_\mathcal{D}>0$ and norm $\lVert .  \rVert_{\mathcal{D}}$ on $\mathbb{R}^d$ such that:
    $$\forall v \in \mathbb{R}^d,\;  \mathbb{E}_{s \sim \mathcal{D} } | \langle v,s \rangle | \ge \mu_{\mathcal{D}} \lVert v\rVert_\mathcal{D}.$$
\end{enumerate}
\end{assumption}
In \cite[Lemma 3.4]{bergou2020stochastic}, the validity of \cref{Ass5} has been established for several distributions including:

 \begin{enumerate}[label=(\roman*)]   
    \item \label{NorDis} For the normal distribution with zero mean and the  identity matrix over $d$ as covariance matrix, i.e., $\mathcal{D} \sim N(0,\frac{I_d}{d})$:
    \[
    \begin{cases}
        \gamma_{\mathcal{D}}=1, \\
        \mathbb{E}_{s \sim \mathcal{D}}|\langle v, s\rangle|=\frac{\sqrt{2}}{\sqrt{d \pi}}\|v\|_2.
    \end{cases}
    \]
     \item \label{unit-sphere}  For the uniform distribution on the unit sphere in \( \mathbb{R}^d \):
$$
    \begin{cases}
    \gamma_{\mathcal{D}}=1,\\
    \mathbb{E}_{s \sim \mathcal{D}} \left| \langle g, s \rangle \right| \sim \frac{1}{\sqrt{2 \pi d}} \|g\|_2.
    \end{cases}
$$
\end{enumerate}

\begin{assumption}\label{Ass6}
 For all $s \sim \mathcal{D}: \mathbb{P}\left( ||s||_2 \le  1 \right)=1.$     
\end{assumption}
Note that under~\Cref{Ass6}, we have $\gamma_{\mathcal{D}}\le 1$. 
Finally, we add the following assumption regarding $\mu_{\mathcal{D}}$ involved in \cref{Ass5}:

\begin{assumption}\label{Ass7}
  $\mu_\mathcal{D} < 1.$  
\end{assumption}

\begin{Remark}\label{Rem1}
    In \Cref{sec5}, we modify the second condition of \Cref{Ass5} by replacing it with:

There exists a constant $\mu_\mathcal{D}>0$ such that: $\forall v \in \mathbb{R}^d,\;  \mathbb{E}_{s \sim \mathcal{D} } | \langle v,s \rangle | \ge \mu_{\mathcal{D}} \lVert v\rVert_2 $.
\end{Remark} 
Since norms are equivalent on $\mathbb{R}^d$, this condition is equivalent to the second condition of \Cref{Ass5}. 
We note also that \Cref{Ass7} is satisfied  for  distribution distribution \ref{NorDis}.
%any arbitrary distribution $\mathcal{D}$ on the set $\{v_1, \ldots, v_n\}$ with probabilities $P(s=v_i)=p_i>0$, and also

Throughout the paper, the abbreviation ``a.s'' stands for ``almost surely''.


\section{Convergence analysis for the class of smooth functions}
\subsection{Convergence analysis for the best iterate}
In this subsection, we will assume that \Cref{Ass1,Ass5,Ass6}  hold true. 
% We then make the assumptions :
% \begin{enumerate}
%      \item $f$ is differentiable, has a lower bound, and is $L$-smooth.
%   \item  The probability distribution $\mathcal{D}$ on $\mathbb{R}^d$ satisfies :
% \begin{itemize}
%     \item $\gamma_{\mathcal{D}}=\mathbb{E}_{s \sim \mathcal{D}}[ \lVert s \rVert_2^2 ]<+\infty$
%     \item There exists a constant $\mu_\mathcal{D}>0$ and norm $\lVert .  \rVert_{\mathcal{D}}$ on $\mathbb{R}^d$ such that:
%     $$(\forall v \in \mathbb{R}^d) :  \mathbb{E}_{s \sim \mathcal{D} } | \langle v,s \rangle | \ge \mu_{\mathcal{D}} \lVert v\rVert_\mathcal{D} $$
% \end{itemize}
% \end{enumerate}
Under these assumptions, we establish
that for any $\epsilon > 0$, when $\{\theta^t\}_{t \ge 1}$ is generated by the STP algorithm using the step size sequence $\{\frac{\alpha}{t^{\frac{1}{2}+\epsilon}} \}_{t \ge 1}$ with $\alpha>0$, it follows that $\min_{1 \le t \le T} \lVert \nabla f(\theta^t) \rVert$ converges almost surely to $0$ at a rate of $o(\frac{1}{T^{\frac{1}{2}-\epsilon }} )$. This result is provided by \Cref{Thm1}, which follows from the first finding of \Cref{Lemma2} that ensures that: $\sum_{t=1}^{\infty} \frac{1}{t^{\frac{1}{2}+\epsilon}} \mathbb{E}\left[ \lVert \nabla f(\theta^t) \rVert_{\mathcal{D}}\right] <\infty.$



\begin{lemma}\label{Lemma2}Assume that \Cref{Ass1,Ass5,Ass6} hold true. Let \(\{\alpha_t\}_{t \ge 1}\) be a sequence of step sizes satisfying \(\sum_{t=1}^{\infty} \alpha_t^2 < \infty\). Let \(\{\theta^t\}_{t \ge 1}\) be a sequence generated by \Cref{Alg1}. Then, the following results hold:
\[
\begin{cases}
 \sum_{t=1}^{\infty} \alpha_t  \mathbb{E}\left[ \lVert \nabla f(\theta^t) \rVert_{\mathcal{D}} \right] < \infty, \\
    \sum_{t=1}^{\infty} \alpha_t \lVert \nabla f(\theta^t) \rVert_{\mathcal{D}} < \infty \quad \text{a.s.}
\end{cases}
\]

\end{lemma}

\vspace{0.1cm}
In the Appendix (Lemma \ref{Lem_Reb}), we prove that if \( \{X_t\}_{t \geq 1} \) is a sequence of nonnegative real numbers that is non-increasing and converges to 0, and \( \{ \alpha_t \}_{t \geq 1} \) is a sequence of real numbers such that \( \sum_{t=1}^{\infty} \alpha_t X_t \) converges, then \( X_T \) converges to 0 at a rate of \( o\left(1 /\sum_{t=1}^T \alpha_t \right) \). As a result, since \( \{\min_{t \leq T} \|\nabla f(\theta^t)\|_{\mathcal{D}} \}_{T \geq 1} \) satisfies the conditions of this lemma when \(\sum_{t=1}^{\infty} \alpha_t^2<\infty \)  and \(\sum_{t=1}^{\infty} \alpha_t = \infty\), we conclude that in this case, the best gradient iterate converges to $0$ at a rate of \( o\left(1 /\sum_{t=1}^T \alpha_t \right) \). This result is formally presented in \Cref{Thm1}.




\begin{theorem}\label{Thm1}
Assume that \Cref{Ass1,Ass5,Ass6} hold. Let \( \{\theta^t\}_{t \ge 1} \) be a sequence generated by \Cref{Alg1}, where the step size sequence \( \{\alpha_t\}_{t \ge 1} \) satisfies the following conditions:
\[
\begin{cases}
    \sum_{t=1}^{\infty} \alpha_t^2 < \infty, \\
    \sum_{t=1}^{\infty} \alpha_t = \infty.
\end{cases}
\]
Then, we have:
\[
\min_{1 \leq t \leq T} \lVert \nabla f(\theta^t) \rVert_{\mathcal{D}} = o\left(\frac{1}{\sum_{t=1}^{T} \alpha_t}\right) \quad \text{a.s.}
\]
In particular, if we choose \( \alpha_t = \frac{\alpha}{t^{\frac{1}{2} + \epsilon}} \) with \( \alpha > 0 \) and \( \epsilon \in \left(0, \frac{1}{2}\right) \), it follows that:
\[
\min_{1 \leq t \leq T} \lVert \nabla f(\theta^t) \rVert_{\mathcal{D}} = o\left(\frac{1}{T^{\frac{1}{2} - \epsilon}}\right) \quad \text{a.s.}
\]

\end{theorem}
In \cite[Theorem 4.1]{bergou2020stochastic}, the authors established that by using the STP algorithm with a step size sequence \(\left\{\frac{\alpha}{\sqrt{t}}\right\}_{t \ge 1}\), where \(\alpha > 0\), the best gradient iterate converges to 0 in expectation at a rate of \(O\left(\frac{\sqrt{d}}{\sqrt{T}}\right)\). The second result of \Cref{Thm1} provides a similar version of this result almost surely, where both the step sizes and convergence rates are roughly similar.

\noindent
\begin{Remark}\label{Rem3}
Since all norms are equivalent in finite dimension, for any norm \( \lVert \cdot \rVert \) on \( \mathbb{R}^d \), we can conclude that by selecting \( \alpha_t = \frac{\alpha}{t^{\frac{1}{2} + \epsilon}} \), where \( \alpha > 0 \) and \( \epsilon \in \left(0, \frac{1}{2}\right) \), the following holds:
\[
\min_{1 \le t \le T} \lVert \nabla f(\theta^t) \rVert = o\left(\frac{1}{T^{\frac{1}{2}-\epsilon}}\right) \quad \text{a.s.}
\]

\end{Remark}
\begin{Remark}\label{Rem4}
In the non-convex setting, the convergence analysis in the previous theorem implies
that $\min_{1 \leq t \leq T} \lVert \nabla f(\theta^t) \rVert$ converges to zero almost surely. However, it remains uncertain whether the gradient of the last iterate $||\nabla f(\theta^T)||$ also converges almost surely to $0$.  In \cref{seclast}, we will establish the convergence of the last iterate of the gradient, both almost surely and in expectation.
\end{Remark}



\subsection{Convergence analysis for the final iterate}\label{seclast}
In this subsection, we will assume that \Cref{Ass1,Ass5,Ass6} hold true. 
% We then make the assumptions 
% \begin{enumerate}
%      \item $f$ is differentiable, has a lower bound, and is $L$-smooth.
%   \item  The probability distribution $\mathcal{D}$ on $\mathbb{R}^d$ satisfies :
% \begin{itemize}
%     \item  For all $s \sim \mathcal{D} : \mathbb{P}\left( ||s||_2 \le  1 \right)=1$
%     \item $\gamma_{\mathcal{D}}=\mathbb{E}_{s \sim \mathcal{D}}[ \lVert s \rVert_2^2 ]<+\infty$
%     \item There exists a constant $\mu_\mathcal{D}>0$ and norm $\lVert .  \rVert_{\mathcal{D}}$ on $\mathbb{R}^d$ such that:
%     $$(\forall v \in \mathbb{R}^d) :  \mathbb{E}_{s \sim \mathcal{D} } | \langle v,s \rangle | \ge \mu_{\mathcal{D}} \lVert v\rVert_\mathcal{D} $$
% \end{itemize}
% \end{enumerate}
Under these assumptions,
we establish that the STP algorithm ensures the almost sure convergence of $||\nabla f(\theta^T)||$ to 0 and the convergence of $\mathbb{E}[ ||\nabla f(\theta^T) ||]$ to 0. This result holds for any step size sequence $\{\alpha_t\}_{t \ge 1}$ such that: $\sum_{t=1}^{\infty} \alpha_t^2 <\infty$ and $\sum_{t =1}^{\infty} \alpha_t=\infty$.
The almost sure convergence result is provided by \Cref{Thm2}, while the convergence in expectation is established by \Cref{Thm3}. Notably, both of these theorems are derived from \Cref{Lemma2} and  \Cref{Lemma3}. (see the Appendix).  

 

\begin{theorem}\label{Thm2}
Assume that \Cref{Ass1,Ass5,Ass6} hold true. Suppose that the step size sequence satisfies:
\[
\begin{cases}
    \sum_{t=1}^{\infty} \alpha_t^2 < \infty, \\
    \sum_{t=1}^{\infty} \alpha_t = \infty.
\end{cases}
\]
Let \( \{\theta^t\}_{t \ge 1} \) be a sequence generated by \Cref{Alg1}. Then, we have:
\[
\lim_{T \to \infty} \lVert \nabla f(\theta^T) \rVert_{\mathcal{D}} = 0 \quad \text{a.s.}
\]

\end{theorem}


\begin{theorem}\label{Thm3}
Assume that \Cref{Ass1,Ass5,Ass6} hold true. Suppose that the step size sequence satisfies:
\[
\begin{cases}
    \sum_{t=1}^{\infty} \alpha_t^2 < \infty, \\
    \sum_{t=1}^{\infty} \alpha_t = \infty.
\end{cases}
\]
Let \( \{\theta^t\}_{t \ge 1} \) be a sequence generated by \Cref{Alg1}. Then, we have:
\[
\lim\limits_{T \rightarrow +\infty} \mathbb{E}[ \lVert \nabla f(\theta^T) \rVert_{\mathcal{D}}]=0 .
\]    
\end{theorem}


\begin{Remark}\label{Rem5}
   In particular, for any $\epsilon \in (0,\frac{1}{2})$, the step size sequence $\{\frac{\alpha}{t^{\frac{1}{2}+\epsilon}} \}_{t \ge 1}$ with $\alpha>0$, satisfies the conditions on step sizes of \Cref{Thm2,Thm3}.
   \end{Remark}
\section{Convergence analysis for the class of smooth convex functions}
In this section we will assume that \Cref{Ass1,Ass2,Ass3,Ass5,Ass6}  hold true. 
Since \( f \) is a real-valued, continuous, and convex function, it follows that \( f \) is a closed proper convex function. Additionally, \Cref{Ass3} guarantees the existence of a vector \( c \) such that the sublevel set of $f$ defined by \( c \) is bounded. Therefore, we can deduce that all sublevel sets of \( f \) are bounded, as shown in \citep[Corollary 8.7.1]{rockafellar2015convex}. 
Let \( \theta^1 \) be the initial vector of the STP algorithm. In particular, the sublevel set $L(\theta^1)$ is bounded, and it forms a compact set of \( \mathbb{R}^d \) (because $f$ is continuous).

Let's denote $\|\cdot\|_{\mathcal{D}}^*$ as the dual norm of $\|\cdot\|_\mathcal{D}$, defined for all $\theta \in \mathbb{R}^d$ by:
$
\|\theta\|_{\mathcal{D}}^* = \sup_{v \in \mathbb{R}^{d} \setminus \{0\}} \frac{\langle v,\theta \rangle}{\|v\|_{\mathcal{D}}}$. 
Since $\theta \mapsto \|\theta-\theta^{*}\|_{\mathcal{D}}^*$ is continuous over $\mathbb{R}^d$ and $L(\theta^1)$ is a compact subset of $\mathbb{R}^d$, we have:  $$R:=\sup_{\theta \in L(\theta^1)} \|\theta-\theta^{*}\|_{\mathcal{D}}^{*} < \infty.$$
Since $f$ is convex, we have that for all $\theta \in L(\theta^1)$:
$$
    f(\theta)-f(\theta^{*}) \le \langle \nabla f(\theta),\theta-\theta^{*} \rangle 
    \le \|\nabla f(\theta)\|_{\mathcal{D}} \underbrace{\sup_{v \in \mathbb{R}^{d} \setminus \{0\}} \frac{\langle v, \theta-\theta^{*} \rangle}{\|v\|_{\mathcal{D}}}}_{\|\theta-\theta^{*} \|_{\mathcal{D}}^*} 
    \le R \|\nabla f(\theta)\|_{\mathcal{D}}.
$$

By the construction of the STP algorithm, for all \( t \ge 1 \), \( f(\theta^t) \le f(\theta^1) \). Therefore, for all \( t \ge 1 \), \( \theta^t \in L(\theta^1) \), and thus we have:
\begin{equation}\label{eq-cvx}
 \forall t \ge 1, \; f(\theta^t)-f(\theta^{*}) \le R \| \nabla f(\theta^t)\|_{\mathcal{D}}.   
\end{equation}

This final result serves as a crucial point for the convergence analysis of \Cref{Thm4} and \Cref{Thm5}.

    

The following~\Cref{Thm4,Thm5}, show the convergence of the final iterate $f(\theta^T)$ to the optimal value with a rate $O(d/T)$ in expectation, and a rate approximately $o(1/T)$ almost surely. 
    

\begin{theorem}\label{Thm4}
Assume that \cref{Ass1,Ass2,Ass3,Ass5,Ass6} hold true, and consider a sequence \( \{\theta^t\}_{t \geq 1} \) generated by \Cref{Alg1}, where the step size sequence is defined as \( \left\{ \frac{\alpha}{t} \right\}_{t \geq 1} \) with \( \alpha > \frac{R}{\mu_{\mathcal{D}}} \).

We have the following bound:
\[
\mathbb{E}\left[ f(\theta^T) \right] - f(\theta^{*}) \leq \frac{a}{T},
\]
where 
\[
a = \max \left( \frac{3\alpha \mu_{\mathcal{D}}}{R} \left(f(\theta^1) - f(\theta^{*})\right), \frac{L \alpha^2}{2 \left( \frac{\alpha \mu_{\mathcal{D}}}{R} - 1 \right)} \right).
\]

In particular, if \( \mu_{\mathcal{D}} \) is proportional to \( \frac{1}{\sqrt{d}} \), then by taking \( \alpha = \frac{2R}{\mu_{\mathcal{D}}} \), we obtain a complexity bound of the form \( O\left( \frac{d}{T} \right) \).



\end{theorem}


Note that for the normal distribution \ref{NorDis} with zero mean and identity covariance matrix , as well as for the uniform distribution over the unit sphere \ref{unit-sphere}, $\mu_{\mathcal{D}}$ is proportional to $\frac{1}{\sqrt{d}}$.

\begin{Remark}\label{LimitR}
Assume that $\mu_{\mathcal{D}}$ is proportional to $\frac{1}{\sqrt{d}}$, and 
let $c\geq 2$ be a constant. If we choose $\alpha$ such that $\frac{2R}{\mu_{\mathcal{D}}}\le \alpha \le \frac{cR}{\mu_{\mathcal{D}}}$, we get $\frac{ L \alpha^2  }{ 2(\frac{\alpha \mu_{\mathcal{D}}  }{R} -1)} = O(c d)$. Thus, we obtain the convergence rate  $O(\frac{c d}{T})$ in \Cref{Thm4}.
%To obtain the convergence rate  $O(\frac{d}{T})$ in \Cref{Thm4}, one must choose the parameter $\alpha$ for the step sizes larger than $\frac{R}{\mu_{\mathcal{D}}}$. However, in practice $R$ can be difficult to determine. %Therefore, one might consider choosing $\alpha$ big enough. 
\end{Remark}


\begin{theorem}\label{Thm5}
Assume that \Cref{Ass1,Ass2,Ass3,Ass5,Ass6} hold true.  Let \( \{\theta^t\}_{t \geq 1} \) be a sequence generated by \Cref{Alg1}, where the step size sequence is given by \( \alpha_t = O\left( \frac{1}{t^{1 - \beta}} \right) \) for some \( \beta \in \left( 0, \frac{1}{2} \right) \).

We then have the following:
\[
\forall \epsilon \in (2\beta, 1), \; f(\theta^T) - f(\theta^{*}) = o\left( \frac{1}{T^{1 - \epsilon}} \right) \quad \text{a.s.}
\]

\end{theorem}




\section{Strongly Convex almost sure convergence rate for STP Algorithm}\label{sec5}
In this section we will assume that \Cref{Ass1,Ass4,Ass5,Ass6,Ass7}  hold true. 
The main results of this section are stated in \Cref{Thm6,Thm7}, which follow from \Cref{Lemma6}.  When step sizes are obtained
by approximating the directional derivatives of the function with respect to the random search directions, 
we show in~\Cref{Thm7} that $f(\theta^T)$ converges in expectation to $\inf_{\theta \in \mathbb{R}^d} f(\theta)$ at a rate of $O((1-\frac{\mu}{2\pi dL})^T)$, 
and in \Cref{Thm6}, we establish this convergence almost surely 
%$f(\theta^T)$ converges almost surely to $\inf_{\theta \in \mathbb{R}^d} f(\theta)$ 
at a rate arbitrarily close to $o((1-\frac{\mu}{2\pi dL})^T)$, 
where $\mu$ and $L$
are  the strong convexity and smoothness parameters of the function, and $d$ is the dimension of the space. 
%where $\beta \in (0,1)$ depends on the strong convexity and smoothness parameters of $f$.
We recall that a strongly convex function has a unique minimizer, which we denote by $\theta^{*}$.% in this section.

The following~\Cref{Lemma5}, controls the decrease per iteration of the value function. It is used to control the total decrease of the value function after $T$ iterations given in~\Cref{Lemma6}. 

\begin{lemma}\label{Lemma5}
Assume that \Cref{Ass1,Ass6} hold true.   Let $h \in (1, \infty)$  and let $\{\theta^t\}_{t \ge 1}$ be a sequence generated by \Cref{Alg1}, where the step size sequence used is $\{ \frac{|f(\theta^t+h^{-t} s_t)-f(\theta^t) |}{L h^{-t}}\}_{t \ge 1}.$ 
Then we have: $$\forall t \ge 1,\;   f(\theta^{t+1}) \le  f(\theta^t)-\frac{  |\langle \nabla f(\theta^t),s_t\rangle|^2 }{2L}+\frac{L}{8}h^{-2t}\; \text{a.s.}$$
\end{lemma}



\begin{lemma}\label{Lemma6}
 Assume that \Cref{Ass1,Ass4,Ass5,Ass6,Ass7}  hold true. Let $h\in (1 , \infty)$ and let $\{\theta^t\}_{t \ge 1}$ be a sequence generated by \Cref{Alg1}, where the step size sequence used is $\{\frac{|f(\theta^t+h^{-t} s_t)-f(\theta^t) |}{L h^{-t}}\}_{t \ge 1}.$ Then we have:
$$\forall T \ge 2,\;
\mathbb{E}[f(\theta^{T})-f(\theta^{*})  ]  \leq\left(1-\frac{\mu_{\mathcal{D}}^2 \mu}{L}\right)^{T-1} [f(\theta^{1})-f(\theta^{*})  ]+ \frac{L}{8}\sum_{i=1}^{T-1}\left(1-\frac{\mu_{\mathcal{D}}^2 \mu}{L}\right)^{T-1-i} h^{-2i}.$$
     
\end{lemma}

\begin{theorem}\label{Thm7}
 Assume that \Cref{Ass1,Ass4,Ass5,Ass6,Ass7} hold true. Let  $h \in \left(\frac{1}{\sqrt{1-\frac{\mu_{\mathcal{D}}^2 \mu}{L}}}, \infty \right)$ and
   let  $\{\theta^t\}_{t \ge 1}$ be a sequence generated by \Cref{Alg1}, where the step size sequence used is $\{\frac{|f(\theta^t+h^{-t} s_t)-f(\theta^t) |}{L h^{-t}}\}_{t \ge 1}$. Then we have:
\begin{equation}\label{eq-strong-exp}
\forall T \ge 2, \;
\mathbb{E}[f(\theta^{T})-f(\theta^{*})  ]  \leq \left(1-\frac{\mu_{\mathcal{D}}^2 \mu}{L}\right)^{T-1} \bigg[f(\theta^{1})-f(\theta^{*})  
+\frac{L}{8}  
\frac{ 1}{h^2 \left(1-\frac{\mu_{\mathcal{D}}^2 \mu}{L} \right)- 1 } \bigg].
\end{equation}

     %\end{itemize}
In particular, if $\mu_{\mathcal{D}}$ is proportional to $\frac{1}{\sqrt{d}}$, i.e., $\mu_{\mathcal{D}} = \frac{K}{\sqrt{d}}$, for some positive constant $K$, then by taking $h=\frac{2}{\sqrt{1-\frac{\mu_{\mathcal{D}}^2 \mu}{L}}}$, we obtain a rate of $O\left(\left(1 - \frac{\mu K^2}{dL}\right)^T\right)$. 
\end{theorem}


\begin{theorem}\label{Thm6}
     Assume that \Cref{Ass1,Ass4,Ass5,Ass6,Ass7} hold true. 
    Let  $\{\theta^t\}_{t \ge 1}$ be a sequence generated by \Cref{Alg1}, where the step size sequence used is $\{\frac{|f(\theta^t+h^{-t} s_t)-f(\theta^t) |}{L h^{-t}}\}_{t \ge 1}$, 
     %\begin{itemize}
      %   \item If \;$\frac{\mu_{\mathcal{D}}^2\mu }{L} =1,$ then : $(\forall T \ge 2) :
    %f(\theta^T) -m=0 \;\; \text{ a.s.}$ (i.e. the STP algorithm converges after one iteration a.s.).
    %\item Otherwise, assuming further that \Cref{Ass7} is satisfied and 
    with $h \in \left(\frac{1}{\sqrt{1-\frac{\mu_{\mathcal{D}}^2 \mu}{L}}}, \infty \right)$, we have: $\forall s \in (0,1),\;  f(\theta^T)-f(\theta^{*}) =o\left((1-  s \frac{\mu_{\mathcal{D}}^2 \mu }{L})^{T}\right)\; \text{ a.s.} $
     %\end{itemize}
     
In particular, if $\mu_{\mathcal{D}}$ is proportional to $\frac{1}{\sqrt{d}}$, i.e., $\mu_{\mathcal{D}} = \frac{K}{\sqrt{d}}$, for some positive constant $K$, then for all $s \in (0,1)$, %by setting $h = \frac{2}{\sqrt{1 - s \frac{\mu_{\mathcal{D}}^2 \mu}{L}}}$, 
we obtain a convergence rate of $o((1- s\frac{\mu K^2}{dL})^T)$. 
\end{theorem}

\begin{Remark}
Note that for the uniform distribution over the unit sphere, we have $\mu_{\mathcal{D}} = \frac{1}{\sqrt{2\pi d}}$. The convergence rates for \Cref{Thm7} and \Cref{Thm6} in this case are respectively $O((1 - \frac{\mu }{2\pi d L })^T)$ and $o((1 - s\frac{\mu }{2 \pi d L})^T)$ for any $s \in (0,1)$. 
\end{Remark}

% \begin{Remark}\label{Rem8}
% In practice, opting for a step-size sequence defined by
% $\{ \frac{|f(\theta^t+h^{-t} s_t)-f(\theta^t) |}{L h^{-t}} \}_{t \ge 1}$
%  may encounter challenges related to machine precision. The primary concern lies in the difficulties of ensuring computational accuracy, as the term 
% $h^{-t}$
%   rapidly diminishes to very small values.
% \end{Remark}


\section{Numerical experiments}
\label{exper}
Let's consider the following optimization problem: 
$$\min_{\theta \in \mathbb{R}^d} f(\theta)=\frac{1}{2}\left( \theta_1\right)^2+\frac{1}{2} \sum_{i=1}^{d-1}\left(\theta_{i+1}-\theta_i\right)^2+\frac{1}{2}\left(\theta_d \right)^2-\theta_1,  \quad \text{ initial vector: } \theta^1=0,$$
where $d=500$. This objective function was used in Section 2.1 of \cite{nesterov2018lectures} to prove the lower complexity bound for gradient methods applied to smooth functions.
%specifically the MNIST model detailed in Table $1$. 
By running multiple trajectories for the three algorithms: the STP algorithm, the RGF algorithm \citep{nesterov2017random}, and the GLD algorithm \citep{golovin2019gradientless}, the objective is to simulate the convergence of the last gradient iterate for each trajectory and also illustrate the rate of convergence of the best gradient iterate. 

\textbf{RGF Algorithm: } This algorithm starts with an
initial vector $\theta^1$
and iteratively updates it according to the following rule
$
\theta^{t+1} = \theta^t - h_t \frac{f(\theta^t + \mu_t u_t) - f(\theta^t)}{\mu_t} u_t,
$
where $u_t$ is a normally distributed gaussian vector. In this implementation, we set $\mu_t = 10^{-4}$ and $h_t = \frac{1}{L}$, where $L\leq 4$ represents the smoothness parameter of the objective function.  %(see \citep[section 8.1]{nesterov2017random})
%The authors of \cite{nesterov2017random} recommend choosing the step size as $h_t = \frac{1}{4L(d+1)}$, where $L$ represents the smoothness parameter of the objective function. For our implementations, we use $h_t = \frac{1}{4 \times 10^{-2} \times (d+1)}$. 

\textbf{GLD algorithm:} This algorithm
proceeds as follows: it starts with an initial point \( \theta^1 \), a sampling distribution \( \mathcal{D} \), and a search radius that shrinks from a maximum value \( R \) to a minimum value \( r \). The number of radius levels is determined by 
\( K = \ceil{\log_2 \left(\frac{R}{r}\right)} \). For each iteration \( t \), the algorithm performs ball sampling trials, where it samples search directions \( v^k \) from progressively smaller radii \( r_k = 2^{-k} R \), $0\leq k \leq K$, and then updates the current point by selecting the \( v^k \) that results in the minimum value of the objective function. The update step is given by:
$
\theta^{t+1} = \argmin_{y\in\{\theta^t,\theta^t+v^0,\cdots,\theta^t+v^K\}}  f(y)$.
For this algorithm, we use the standard Gaussian distribution $\mathcal{D}$ and set $r=10^{-5}$ and $R=
10^{-4}$.


For the STP algorithm, we set the step sizes to be \(\alpha_t = \frac{4}{t^{0.51}}\), and the random search directions \(s_t\) are generated uniformly on the unit sphere of \(\mathbb{R}^d\). The chosen step sizes adhere to the form provided in the second result of~\Cref{Thm1}, where \(\epsilon = 0.01\). In our experiment, we run 50 trajectories for each of the three algorithms, all starting from the same initial point \(0\). We simulate \(\log_{10}(\|\nabla f(\theta^T)\|_2)\) as a function of the number of iterations, as well as the elapsed time in seconds. Additionally, to verify the rate assured by~\Cref{Thm1} for the STP algorithm, we simulate $T^{0.49}\min_{t \le T} \|\nabla f(\theta^t)\|_2$ as a function of the number of iterations. 

\Cref{fig:lastiter} and \Cref{fig:vstime} illustrate the logarithmic decay of the gradient norm with respect to both iterations and elapsed time, highlighting its convergence to zero across all trajectories for the three algorithms. Notably, STP and RGF demonstrate competitive performance, with STP being slightly better, in terms of the number of iterations and the time required to achieve a given accuracy, outperforming the GLD method in both metrics. This similarity between the performance of STP and RGF reflects their similar theoretical complexity bounds. It is important to note also that at each iteration, the STP and RGF methods require two function evaluations, while the GLD method requires $\ceil{\log_2 \left(\frac{R}{r}\right)}$ function evaluations.

In \Cref{fig:rateiter}, we observe the convergence of the best gradient iterate to $0$ at a rate of $o(\frac{1}{T^{0.49}})$ across all trajectories for the three algorithms. In particular, this illustrates the rate obtained for the STP algorithm.


\begin{figure}[H]
    \centering
    \includegraphics[width=1\linewidth]{Last_iteration_log.png}
    \caption{Logarithmic decay of gradient norm vs. Iterations.}
    \label{fig:lastiter}
\end{figure}


\begin{figure}[H]
    \centering
    \includegraphics[width=1\linewidth]{time.png}
    \caption{Logarithmic Decay of gradient norm vs. Time.}
    \label{fig:vstime}
\end{figure}

\begin{figure}[H]
    \centering
    \includegraphics[width=1\linewidth]{Best_iteration_log.png}
    \caption{Convergence rate of the best gradient iterate.}
    \label{fig:rateiter}
\end{figure}
\color{black}


\bibliography{iclr2025_conference}
\bibliographystyle{iclr2025_conference}

\appendix
\section{Appendix}
\begin{lemma}\label{Lemma1} \textbf{\cite[Lemma 3.5]{bergou2020stochastic}}
Assume that \Cref{Ass1,Ass5,Ass6} hold true and let $\{\theta^t\}_{t \ge 1}$ be a sequence generated by \Cref{Alg1}. We have:
    $$
    \mathbb{E}[f(\theta^{t+1})  \: \mid \; \theta^t  ]
   \le f(\theta^t) - \mu_{\mathcal{D}} \alpha_t \lVert \nabla f(\theta^t) \rVert_{\mathcal{D}} + \frac{L\alpha_t^2 }{2}.$$

\end{lemma}
\vspace{0.1cm}

\begin{proof}[Proof of \Cref{Lemma2}]

Let $t \ge 1$. By~\Cref{Lemma1}, we have that: 
$$
    \mathbb{E}[f(\theta^{t+1})  \: \mid \; \theta^t  ]
   \le f(\theta^t) - \alpha_t\mu_{\mathcal{D}} \lVert \nabla f(\theta^t) \rVert_{\mathcal{D}} + \frac{L\alpha_t^2  }{2}.
$$ 
By taking the expectation, we get:
$
    \mathbb{E}[f(\theta^{t+1}) ]
   \le \mathbb{E}[f(\theta^t)] - \alpha_t\mu_{\mathcal{D}} \mathbb{E}[\lVert \nabla f(\theta^t) \rVert_{\mathcal{D}}] + \frac{L\alpha_t^2  }{2}.
$ 

It follows that: 
\begin{equation}\label{eq-1}
 \mu_{\mathcal{D}}  \alpha_t \mathbb{E}[\lVert \nabla f(\theta^t) \rVert_\mathcal{D}] \le \mathbb{E}[f(\theta^{t}) ]-\mathbb{E}[f(\theta^{t+1}) ]  + \frac{L\alpha_t^2  }{2}. 
\end{equation}

By construction of the algorithm the sequence $\{f(\theta^{t})\}_{t\ge 1}$ is non-increasing, and since we assume that $f$ is bounded from below, we have that 
$\{\mathbb{E}[f(\theta^{t})]\}_{t \ge  1}$ is non-increasing and bounded from bellow, and thus converges. As a result, we have: 
$\sum_{t=1}^{\infty} \mathbb{E}[f(\theta^{t}) ]-\mathbb{E}[f(\theta^{t+1}) ]<\infty$. 
Knowing that $\sum_{t=1}^{\infty}\alpha_t^2 < \infty$, we conclude from~\eqref{eq-1}, that 
$$ \sum_{t=1}^{\infty} \alpha_t  \mathbb{E}\left[ \lVert \nabla f(\theta^t) \rVert_{\mathcal{D}}\right] <\infty.$$


We deduce also that 
$
\mathbb{E}\left[\sum_{t=1}^{\infty} \alpha_t \lVert \nabla f(\theta^t) \rVert_\mathcal{D} \right]  =\sum_{t=1}^{\infty} \alpha_t \mathbb{E}\left[ \lVert \nabla f(\theta^t) \rVert_{\mathcal{D}}\right] <\infty,$ which implies that:
$$
\sum_{t=1}^{\infty} \alpha_t \lVert \nabla f(\theta^t) \rVert_\mathcal{D} < \infty \text{ a.s.} 
$$
\end{proof}

\begin{lemma}\label{Lem_Reb} 
 Let \( \{X_t\}_{t \ge 1} \) be a sequence of nonnegative real numbers that is non increasing and converges to \( 0 \), and let \( \{\alpha_t\}_{t \ge 1} \) be a sequence of real numbers such that  \( \sum_{t=1}^{\infty} \alpha_t X_t \) converges. Then, we have:
\[
X_T = o\left( \frac{1}{\sum_{t=1}^T \alpha_t} \right).
\]


\end{lemma}
\begin{proof}
    For all $T \ge 1$, we define $U_T=X_T \sum_{i=1}^T \alpha_i$ and $R_T=\sum_{i=T}^{\infty} \alpha_i X_i.$ We then have:
 
 $$
    U_T=X_T\sum_{i=1}^T (R_{i}-R_{i+1})\frac{1}{X_i} .
  $$
 Let $T \ge 2.$ We have:
 \begin{align*}
      U_T&=X_T \left[ \sum_{i=1}^T R_i\frac{1}{X_i}-\sum_{i=1}^T R_{i+1}\frac{1}{X_i}  \right]\\
      &= X_T \left[ \sum_{i=1}^T R_i\frac{1}{X_i}-\sum_{i=2}^{T+1} R_i\frac{1}{X_{i-1}}   \right]\\
      &= X_T \left[ R_1 \frac{1}{X_1}-\frac{R_{T+1}}{X_{T}} +\sum_{i=2}^T R_i (\frac{1}{X_i}-\frac{1}{X_{i-1}})  \right]\\
      &=R_1 \frac{X_T}{X_1}-R_{T+1}+X_T \sum_{i=2}^T R_i (\frac{1}{X_i}-\frac{1}{X_{i-1}}).
 \end{align*}
 To prove \( \lim_{T \to +\infty} U_T = 0 \), it suffices to show that:
\[
\lim_{T \to +\infty} X_T \sum_{i=2}^T R_i (\frac{1}{X_i}-\frac{1}{X_{i-1}}) = 0.
\]

Let $\epsilon>0$ and $T_0 \ge 2$ such that for all $T \ge T_0$, we have $ R_T  \le \frac{\epsilon}{2}$. Let $T > T_0$, we have:
 \begin{align*}
 \left| X_T \sum_{i=2}^T R_i (\frac{1}{X_i}-\frac{1}{X_{i-1}})    \right| &\le X_T \sum_{i=2}^{T_0} |R_i|(\frac{1}{X_i}-\frac{1}{X_{i-1}}) +X_T\sum_{i=T_0+1}^{T} \frac{\epsilon}{2} (\frac{1}{X_i}-\frac{1}{X_{i-1}})
 \\
 &= X_T \sum_{i=2}^{T_0} |R_i|(\frac{1}{X_i}-\frac{1}{X_{i-1}}) + \frac{\epsilon X_T}{2}  (\frac{1}{X_T}-\frac{1}{X_{T_0}})
  \\
 &= X_T \sum_{i=2}^{T_0} |R_i|(\frac{1}{X_i}-\frac{1}{X_{i-1}}) + \frac{\epsilon}{2} (1-\frac{X_T}{X_{T_0}})
 \\&\le X_T \sum_{i=2}^{T_0} |R_i|(\frac{1}{X_i}-\frac{1}{X_{i-1}}) + \frac{\epsilon}{2}.
 \end{align*}
As $\lim_{T \to +\infty} X_T=0$, there exists $T_1 \ge T_0$ such that for all $T \ge T_1$,
 \begin{align*}
 \left| X_T \sum_{i=2}^T R_i (\frac{1}{X_i}-\frac{1}{X_{i-1}})    \right| 
 &\le X_T \sum_{i=2}^{T_0} |R_i|(\frac{1}{X_i}-\frac{1}{X_{i-1}}) + \frac{\epsilon }{2}\le \epsilon.
 \end{align*}
 Therefore $\lim_{T \to +\infty} X_T \sum_{i=2}^T R_i (\frac{1}{X_i}-\frac{1}{X_{i-1}}) = 0,$ and we deduce that: \[
X_T = o\left( \frac{1}{\sum_{t=1}^T \alpha_t} \right) \text{ as } T \to +\infty.
\]
\end{proof}


The following lemma, which is a classical result about Riemann series, will be needed in the proof of~\Cref{Thm1}.
\begin{lemma}\label{Lemma8}
    For all $\alpha \in (0,1)$, we have: $\sum_{t=1}^T \frac{1}{t^\alpha}\sim \frac{T^{1-\alpha}}{1-\alpha}$. 
\end{lemma}
%\textbf{Proof of~\Cref{Thm1}.}
\begin{proof}[Proof of \Cref{Thm1}]
Let us define \( X_T = \min_{t \le T} \|\nabla f(\theta_t) \|_{\mathcal{D}} \) for all \( T \geq 1 \). Since $\sum_{t= 1}^{\infty} \alpha_t^2<\infty $, according to ~\Cref{Lemma2}, we deduce that \( \sum_{t=1}^{\infty} \alpha_t X_t < \infty \) a.s. It is clear that $\{X_T\}_{T \ge 1}$ is a sequence of nonnegative real numbers that is non increasing, then, by proving $\lim_{T \to +\infty} X_T =0 \text{ a.s.}$, using \cref{Lem_Reb} , we can deduce that:

\[
X_T = o\left(\frac{1}{\sum_{t=1}^{T} \alpha_t}\right) \text{a.s.} 
\]
Now, we prove that \( \lim\limits_{T \to \infty} X_T = 0 \) a.s. According to ~\Cref{Lemma2}, we have \( \sum_{t=1}^{\infty} \alpha_t ||\nabla f(\theta^t)||_{\mathcal{D}} < \infty \) a.s. Thus, it follows that:
\[
\{ (\min_{t \leq T} \|\nabla f(\theta^t)\|) \sum_{t=1}^T \alpha_t\}  \quad \text{is bounded almost surely}.
\]
Since $\lim\limits_{T \to +\infty} \sum_{t=1}^T \alpha_t = +\infty$, we can conclude that:
\[
\lim_{T \to +\infty} X_T = \lim_{T \to +\infty} \min_{t \leq T} \|\nabla f(\theta^t)\|_{\mathcal{D}} = 0 \text{ 
 a.s.} 
\]

Therefore, we establish the first result of the theorem. The second result is obtained by choosing \( \{\alpha_t\}_{t \ge 1} \) defined by \( \alpha_t = \frac{1}{t^{\frac{1}{2} + \epsilon}} \). In this case, we have \( \sum_{t=1}^{\infty} \alpha_t = \infty \), while \( \sum_{t=1}^{\infty} \alpha_t^2 <\infty\).


Using \Cref{Lemma8}, we have \[
\sum_{t=1}^{T} \frac{1}{t^{\frac{1}{2}+\epsilon}} \sim  \frac{T^{\frac{1}{2}-\epsilon}}{\frac{1}{2}-\epsilon} .
\]

Therefore, \[
\min_{1 \leq t \leq T} \lVert \nabla f(\theta^t) \rVert_{\mathcal{D}} = o\left(\frac{1}{T^{\frac{1}{2} - \epsilon}}\right) \quad \text{a.s.} 
\]
\end{proof}



\Cref{Lemma3} is first presented in \cite[Proposition 2]{alber1998projected} and again in \cite[Lemma A.5]{mairal2013stochastic}, along with a new proof. We provide a new, simpler proof of this lemma that is more straightforward than those presented in these references.
\begin{lemma}\label{Lemma3}
 \textbf{ \cite[Proposition 2]{alber1998projected} , \cite[Lemma A.5]{mairal2013stochastic} } Let $\{a_t\}_{t \geq 1},\{b_t\}_{t \geq 1}$ be two nonnegative real sequences. We have:
$$
\begin{cases}
     \sum_{t=1}^{\infty} a_t b_t<\infty,\\
     \sum_{t=1}^{\infty} a_t=\infty,\\
     \text{There exists } K \geq 0 \text{ such that } \left|b_{t+1}-b_t\right| \leq K a_t. 
 \end{cases} \Longrightarrow
\lim\limits_{t \rightarrow +\infty} b_t = 0.
$$
\end{lemma}
\noindent


\begin{proof}[Proof of \Cref{Lemma3}]
First, we note that for all \( n_0 \geq 1 \), we have \( \inf_{n \geq n_0} b_n = 0 \). Indeed, suppose for contradiction that \( \inf_{n \geq n_0} b_n > 0 \). In this case, we would have for all $n \ge n_0$, \( a_n b_n \geq a_n \inf_{m \geq n_0} b_m \), which implies that the series \( \sum a_n b_n \) cannot converge, since \( \sum_{n=1}^{\infty} a_n=\infty \)  and $inf_{m \ge n_0} b_m>0$. This contradiction implies that for all \( n_0 \geq 1 \), we have \( \inf_{n \geq n_0} b_n = 0 \).


Let $\epsilon>0$. Let $n_0 \ge 1$  such that for all $n \ge n_0$ we have $\sum_{k=n}^{\infty} a_k b_k \le \frac{\epsilon^2}{4K}.$ 

The goal is to prove that for all \( n \geq n_0 \), \( b_n \leq \epsilon \). Let \( n \geq n_0 \). If \( b_n \leq \frac{\epsilon}{2} \), then trivially \( b_n \leq \epsilon \). Now assume that \( b_n > \frac{\epsilon}{2} \).

We have \( \inf_{t \geq n} b_t = 0 \), then we can take the smallest index \( m > n \) such that \( b_m \leq \frac{\epsilon}{2} \). We have: 
\begin{align*}
    |b_m-b_n| &\le \sum_{i=n}^{m-1} |b_{i+1}-b_i| \\
    &\le K \sum_{i=n}^{m-1} a_i \\
    &= K \sum_{i \in \{n,...,m-1\},b_i > \frac{\epsilon}{2}} a_i \\
    &\le \frac{2K}{\epsilon}\sum_{i \in \{n,...,m-1\},b_i > \frac{\epsilon}{2}} a_i b_i\\
    &\le \frac{2K}{\epsilon} \sum_{i=n}^{\infty} a_i b_i
    \\
    &\le \frac{\epsilon}{2}.
\end{align*}
Therefore, by the triangle inequality, we have:
\begin{align*}
b_n &\le b_m+\frac{\epsilon}{2} \le \epsilon.
\end{align*}
Thus, for all \( n \geq n_0 \), we have \( b_n \leq \epsilon \), and consequently, we deduce that \( \lim_{n \to +\infty} b_n = 0 \).




\end{proof}


%\textbf{Proof of~\Cref{Thm2}.}
\begin{proof}[Proof of \Cref{Thm2}]

Consider $C > 0$ satisfying: $
||.||_{\mathcal{D} } \le C ||.||_2$.\\\\
Let $t\ge 1.$ We have that:
\begin{align*}
    \left|\left\|\nabla f\left(\theta^{t+1}\right)\right\|_{\mathcal{D}}-\left\|\nabla f\left(\theta^t\right)\right\|_{\mathcal{D}}\right| &\le \lVert \nabla f(\theta^{t+1})-\nabla f(\theta^t) \rVert_{\mathcal{D}} \\
    &\le C L \lVert \theta^{t+1}-\theta^t \rVert_{2} \text{ (because } f \text{ is } L\text{-smooth)}\\
    &= CL\alpha_t \lVert s_t \rVert_{2}\\
    &\le CL\alpha_t \;\; \text{ a.s.}  \text{  (because we assume~\Cref{Ass6} holds true)}
\end{align*}

Therefore, we have that for all $t \ge 1$: $\mathbb{P}\left( 
 \left|\left\|\nabla f\left(\theta^{t+1}\right)\right\|_{\mathcal{D}}-\left\|\nabla f\left(\theta^t\right)\right\|_{\mathcal{D}}\right|  \le CL \alpha_t \right)=1.$ 
 
 Thus: $\mathbb{P}\left( \forall t \ge 1,\;  
 \left|\left\|\nabla f\left(\theta^{t+1}\right)\right\|_{\mathcal{D}}-\left\|\nabla f\left(\theta^t\right)\right\|_{\mathcal{D}}\right|  \le CL \alpha_t \right)=1.$

Given that
$\begin{cases}
\sum_{t=1}^{\infty} \alpha_t   \lVert \nabla f(\theta^t) \rVert_{\mathcal{D}}< \infty \text { a.s, \; (by~\Cref{Lemma2} because $\sum_{t=1}^{ \infty} \alpha_t^2 <\infty$)}\\
    \sum_{t=1}^{\infty} \alpha_t=\infty .
\end{cases}$

Using~\Cref{Lemma3}, with $\{\alpha_t\}_{t\ge 1}$ playing the role of $\{a_t\}_{t\ge 1}$ and $\{\lVert \nabla f(\theta^t) \rVert\}_{t\ge 1}$ playing the role of $\{b_t\}_{t\ge 1}$, we conclude that: $$\lim\limits_{T \rightarrow +\infty} \lVert \nabla f(\theta^T) \rVert_{\mathcal{D}}=0\;\; \text{ a.s.}$$ 
\end{proof}

%\textbf{Proof of~\Cref{Thm3}.} 
\begin{proof}[Proof of \Cref{Thm3}]
Consider $C > 0$ satisfying: $
||.||_{\mathcal{D} } \le C ||.||_2$.\\\\
Let $t\ge 1$. We have that:
\begin{align*}
    \left| \mathbb{E}\left[ \left\|\nabla f\left(\theta^{t+1}\right)\right\|_{\mathcal{D}}\right]-\mathbb{E}\left[\left\|\nabla f\left(\theta^t\right)\right\|_{\mathcal{D}}\right] \;\right|  
    &\le
   \mathbb{E}\left[ \; \Big|  \left\|\nabla f\left(\theta^t\right)\right\|_{\mathcal{D}}-  \left\|\nabla f\left(\theta^{t+1}\right)\right\|_{\mathcal{D}}  \Big| \;\right]  
    \\
     &\le
   \mathbb{E}\left[ \lVert \nabla f\left(\theta^{t}\right)-\nabla f\left(\theta^{t+1}\right)  \rVert_{\mathcal{D}} \right]\\
    &\le CL \mathbb{E}\left[ 
  \lVert\theta^{t+1}-\theta^t\rVert_2 \right] \text{ (because } f \text{ is } L\text{-smooth)}
  \\ &= CL \alpha_t \mathbb{E}\left[ \lVert s_t \rVert_2 \right]
  \\ &\le CL \alpha_t  \text{  (because we assume~\Cref{Ass6} holds true)},
\end{align*}
where in the first inequality, we used Jensen's inequality. So, we proved that:
%It holds that : 
$$\forall t \ge 1,\; 
 \left| \mathbb{E}[\left\|\nabla f\left(\theta^{t+1}\right)\right\|_{\mathcal{D}}]-\mathbb{E}[\left\|\nabla f\left(\theta^t\right)\right\|_{\mathcal{D}}]\right|  \le CL \alpha_t .$$


Now, given that
$\sum_{t=1}^{\infty} \alpha_t   \mathbb{E}[ \lVert \nabla f(\theta^t) \rVert_{\mathcal{D}}] <\infty$  (by~\Cref{Lemma2} because $\sum_{t=1}^{ \infty} \alpha_t^2 <\infty$), and that
$\sum_{t=1}^{\infty} \alpha_t=\infty$ , we use~\Cref{Lemma3}, with $\{\alpha_t\}_{t\ge 1}$ playing the role of $\{a_t\}_{t\ge 1}$ and $\{\mathbb{E}[\lVert \nabla f(\theta^t) \rVert]\}_{t\ge 1}$ playing the role of $\{b_t\}_{t\ge 1}$, to conclude that
$$\lim\limits_{T \rightarrow +\infty} \mathbb{E}[ \lVert \nabla f(\theta^T) \rVert_{\mathcal{D}}]=0.$$
\end{proof}



\begin{lemma}\label{Lemma4}\textbf{\cite[Lemma 1]{liu2022almost} }
 If $\left\{Y_t\right\}_{t \ge 1}$ is a sequence of nonnegative random variables adapted to a filtration $\left\{\mathcal{F}_t\right\}_{t \ge 1}$, and satisfying:
$$
\mathbb{E}\left[Y_{t+1} \mid \mathcal{F}_t\right] \leq\left(1-c_1 \alpha_t\right) Y_t+c_2 \alpha_t^2\;\; \text{ for all } \;\;t \geq 1 ,
$$
where $\alpha_t=O\left(\frac{1}{t^{1-\beta}}\right)$ for some $\beta \in\left(0, \frac{1}{2}\right)$, and $c_1$ and $c_2$ are positive constants. Then, for any $\epsilon \in(2 \beta, 1):$
$$
Y_t=o\left(\frac{1}{t^{1-\epsilon}}\right) \quad \text { a.s.} 
$$
\end{lemma}



%\textbf{Proof of~\Cref{Thm4}.}
\begin{proof}[Proof of \Cref{Thm4}]
By~\Cref{Lemma1}, we have that:
$$
    \forall t \ge 1,\; \mathbb{E}[f(\theta^{t+1})  \: \mid \; \theta^t  ]
   \le f(\theta^t) - \mu_{\mathcal{D}} \alpha_t \lVert \nabla f(\theta^t) \rVert_{\mathcal{D}} + \frac{L \alpha_t^2 }{2}.
$$
Knowing from~\eqref{eq-cvx} that: $\forall t \ge 1,\;  f(\theta^t)-f(\theta^{*}) \le R \| \nabla f(\theta^t)\|_{\mathcal{D}}$, we have:
$$\forall t \ge 1,\;  \mathbb{E}[f(\theta^{t+1})-f(\theta^{*}) \: \mid \; \theta^t  ] \le \left( 1- \alpha_t \frac{\mu_{\mathcal{D}}}{ R}  \right) \left(f(\theta^{t})-f(\theta^{*})\right)+ \frac{L  \alpha_t^2  }{2}.$$
Taking the expectation, and knowing that $\alpha_t=\frac{\alpha}{t}$, we get: 
\begin{equation}\label{eq-delta-cvx}
  \forall t \ge 1,\;   \underbrace{\mathbb{E}\left[f(\theta^{t+1})-f(\theta^{*})  \right]}_{:=\delta_{t+1}} \le \left( 1-  \frac{\alpha \mu_{\mathcal{D}}}{ R t}  \right) \underbrace{\mathbb{E}\left[f(\theta^{t})-f(\theta^{*})\right]}_{:=\delta_t}+ \frac{L  \alpha^2  }{2 t^2}.
\end{equation}
If $t \in \{1,...,[\frac{\alpha \mu_{\mathcal{D}}}{R} ]+1  \}$, $\delta_t=\mathbb{E}\left[f(\theta^{t})-f(\theta^{*})\right] \le f(\theta^1)-f(\theta^{*}) \le ([\frac{\alpha \mu_{\mathcal{D}}}{R} ]+1) \frac{ f(\theta^1)-f(\theta^{*})}{t}$. Then $$\forall t \in \{1,...,[\frac{\alpha \mu_{\mathcal{D}}}{R} ]+1  \},\; \delta_t \le \frac{3\alpha \mu_{\mathcal{D}} }{R}\frac{f(\theta^1)-f(\theta^{*})}{t}\;\; \text{, because } [\frac{\alpha \mu_{\mathcal{D}} }{R}] +1 \le  \frac{\alpha \mu_{\mathcal{D}} }{R}+2 \le \frac{3\alpha \mu_{\mathcal{D}} }{R}.$$

Let's denote $a$ and $b$ as follows:  $
a= \max \left(  \frac{3\alpha \mu_{\mathcal{D}} }{R}(f(\theta^1)-f(\theta^{*})), \frac{ L  \alpha^2  }{ 2(\frac{\alpha \mu_{\mathcal{D}}  }{R} -1)} \right)
$ and $
b=\frac{\mu_{\mathcal{D}}}{R}$. 
We have:
\begin{equation}\label{Eqrec}
    \forall  t  \in \{1,...,[\frac{\alpha \mu_{\mathcal{D}}}{R} ]+1  \}, \;  \delta_t \le \frac{ a}{t }.
\end{equation}

We will prove by induction that: 
\begin{equation*}
    \forall t \ge [\frac{\alpha \mu_{\mathcal{D}} }{R}]+1,\;   \delta_t \le \frac{ a}{t }.
\end{equation*}
For $t=[\frac{\alpha \mu_{\mathcal{D}} }{R}]+1$, we have that: $$\delta_{[\frac{\alpha \mu_{\mathcal{D}} }{R}]+1} \le  \frac{a}{t}.$$

Let $t \ge [\frac{\alpha \mu_{\mathcal{D}} }{R}]+1$.  Assume that $\delta_t \le \frac{ a}{t }$ and let's prove that  $\delta_{t+1}\le \frac{a}{t+1}$. We note that $1-  \frac{\alpha \mu_{\mathcal{D}}}{ R t}>0$. 

From~\eqref{eq-delta-cvx}, we get 
$
 \delta_{t} \le \frac{a}{t} \Longrightarrow   \delta_{t+1} \le \frac{a}{t}-\frac{ab\alpha}{t^2}+\frac{L \alpha^2 }{2 t^2}$. 
We have also the following equivalence:
\begin{align*}
\frac{a}{t}-\frac{ab\alpha}{t^2}+\frac{L  \alpha^2}{2  t^2}   \le \frac{a}{t+1} &\Longleftrightarrow -\frac{a b\alpha}{t^2}+\frac{L \alpha^2 }{2  t^2}  \le \frac{-a}{t(t+1)}
\\
&\Longleftrightarrow -ab\alpha +\frac{L \alpha^2}{2 } \le \frac{-at}{t+1}\enspace.
\end{align*}

Let's prove that the last assertion is true. We have: 
\begin{align*}
    -a \le \frac{-at}{t+1} &\Longrightarrow -ab\alpha+a(b \alpha-1)\le  \frac{-at}{t+1}\\
    &\Longrightarrow -ab\alpha+\frac{L \alpha^2}{2 }\le  \frac{-at}{t+1}\;.
\end{align*}
The last implication comes from $a(b \alpha-1)=\max\left((\underbrace{b\alpha-1}_{>0})  \frac{3\alpha \mu_{\mathcal{D}} }{R}(f(\theta^1)-f(\theta^{*})), \frac{L \alpha^2}{2 }   \right)$.

We deduce finally that $\delta_{t+1} \le \frac{a}{t+1}$. Therefore, we get: $\forall T \ge [\frac{\alpha \mu_{\mathcal{D}} }{R}]+1,\;  \mathbb{E}[f(\theta^T)]-f(\theta^{*}) \le \frac{a}{T}$, 
and using \eqref{Eqrec}, we deduce that: $$\forall T \ge 1,\;  \mathbb{E}[f(\theta^T)]-f(\theta^{*}) \le \frac{a}{T}.$$

In particular, if $\mu_{\mathcal{D}}$ is proportional to $\frac{1}{\sqrt{d}}$, then by taking  $\alpha=\frac{2R}{\mu_{\mathcal{D}}}$, we have: $$a=\max( \frac{3\alpha \mu_{\mathcal{D}} }{R}(f(\theta^1)-f(\theta^{*})), \frac{L \alpha^2 }{2( \frac{\alpha \mu_{ \mathcal{D}} }{R}-1 ) })=\max ( 6(f(\theta^1)-f(\theta^{*})),   \frac{L\frac{4R^2}{\mu_{\mathcal{D}}^2} }{2})=O(d),$$
therefore $\mathbb{E}[f(\theta^T)]-f(\theta^{*})=O(\frac{d}{T})$.
\end{proof}

%\textbf{Proof of~\Cref{Thm5}.}
\begin{proof}[Proof of \Cref{Thm5}]

By employing the first part of the proof of~\Cref{Thm4}, we have that:
$$\forall t \ge 1,\;  \mathbb{E}[f(\theta^{t+1})-f(\theta^{*}) \: \mid \; \theta^t  ] \le \left( 1- \alpha_t \frac{\mu_{\mathcal{D}}}{ R}  \right) \left(f(\theta^{t})-f(\theta^{*})\right)+ \frac{L \alpha_t^2  }{2}$$
Using~\Cref{Lemma4}, we deduce that when $\alpha_t=O(\frac{1}{t^{1-\theta}})$ with $\theta \in (0,\frac{1}{2})$, we get: 
$$\forall \epsilon \in (2\theta,1),\; f(\theta^T)-f(\theta^{*})=o(\frac{1}{T^{1-\epsilon}}) \;\;\text{ a.s.}$$
\end{proof}


\begin{proof}[Proof of \Cref{Lemma5}]
%\textbf{Proof of~\Cref{Lemma6}.}

Let $t \ge 1$.
Using the smoothness property in~\eqref{ineq-smooth}, we have: 
$$
\begin{cases}
f(\theta^{t}+ \alpha_t s_t ) \le f(\theta^t)+ \alpha_t \langle \nabla f(\theta^t),s_t\rangle+\frac{L}{2} \alpha_t^2 ||s_t||^2 \\
f(\theta^{t}- \alpha_t s_t ) \le f(\theta^t)- \alpha_t \langle \nabla f(\theta^t),s_t\rangle+\frac{L}{2} \alpha_t^2 ||s_t||^2
\end{cases} .
$$

Then $f(\theta^{t+1}) \le f(\theta^t)- \alpha_t | \langle \nabla f(\theta^t),s_t \rangle| +\frac{L}{2} \alpha_t^2 ||s_t||^2$. By replacing $\alpha_t$ by its expression, and using~\Cref{Ass6}, we have: 
\begin{align*}
      f(\theta^{t+1}) &\le f(\theta^t)- \frac{|f(\theta^t+h^{-t} s_t)-f(\theta^t) |}{L h^{-t}} | \langle \nabla f(\theta^t),s_t \rangle|+\frac{L}{2} \left(\frac{f(\theta^t+h^{-t} s_t)-f(\theta^t) }{L h^{-t}} \right)^2 \text{ a.s.} \\
       &\le f(\theta^t)- \frac{|f(\theta^t+h^{-t} s_t)-f(\theta^t) |}{L h^{-t}} | \langle \nabla f(\theta^t),s_t \rangle|\\&+\frac{L}{2} \left( \frac{\lvert f(\theta^t+h^{-t}s_t )-f(\theta^t)\rvert- \lvert\langle \nabla f(\theta^t),h^{-t}s_t \rangle \rvert }{Lh^{-t}}  \right)^2\\&+\frac{ |f(\theta^t+h^{-t}s_t )-f(\theta^t)|\; | \langle \nabla f(\theta^t),h^{-t}s_t \rangle| }{Lh^{-2t}}-\frac{  |\langle \nabla f(\theta^t),h^{-t}s_t\rangle|^2 }{2L h^{-2t}}\text{ a.s.}
\\
      &\le f(\theta^t)-\frac{  |\langle \nabla f(\theta^t),s_t\rangle|^2 }{2L}+\frac{L}{2} \left( \frac{|f(\theta^t+h^{-t}s_t )-f(\theta^t)|- |\langle \nabla f(\theta^t),h^{-t}s_t \rangle|  }{Lh^{-t}}  \right)^2\text{ a.s.}
      \\
      &\le f(\theta^t)-\frac{  |\langle \nabla f(\theta^t),s_t\rangle|^2 }{2L}+\frac{L}{2} \left( \frac{|f(\theta^t+h^{-t}s_t )-f(\theta^t)- \langle \nabla f(\theta^t),h^{-t}s_t \rangle | }{Lh^{-t}}  \right)^2 \text{ a.s.}
            \\
      &\le f(\theta^t)-\frac{  |\langle \nabla f(\theta^t),s_t\rangle|^2 }{2L}+\frac{L}{2} \left( \frac{ \frac{L h^{-2t} ||s_t||^2 }{2}   }{Lh^{-t}}  \right)^2 \text{ a.s.  }  (\text{using property~\eqref{ineq-smooth}} )
      \\ &\le  f(\theta^t)-\frac{  |\langle \nabla f(\theta^t),s_t\rangle|^2 }{2L}+\frac{L}{8}h^{-2t}\;\; \text{ a.s.  }
    \end{align*}
  We conclude that: $\forall t \ge 1,\;  f(\theta^{t+1}) \le  f(\theta^t)-\frac{  |\langle \nabla f(\theta^t),s_t\rangle|^2 }{2L}+\frac{L}{8}h^{-2t}\; \text{a.s.} $ 
\end{proof}

%\textbf{Proof of~\Cref{Lemma7}.}
\begin{proof}[Proof of \Cref{Lemma6}]

Let \( t \ge 1 \). By~\Cref{Lemma5}, we have: 
\[ 
f(\theta^{t+1}) \leq f(\theta^t) - \frac{\left|\left\langle \nabla f(\theta^t), s_t \right\rangle\right|^2}{2L} + \frac{Lh^{-2t}}{8} \quad \text{a.s.} 
\]


Using the tower property: 
\begin{align*}
\mathbb{E}\left[\left|\left\langle\nabla f\left(\theta^t \right), s_t\right\rangle\right|^2\right] 
&= \mathbb{E}\left[\mathbb{E}_{s_t\sim \mathcal{D}}\left[\left|\left\langle\nabla f\left(\theta^t\right), s_t\right\rangle\right|^2 \mid \theta^t\right]\right] \\
&\underbrace{\geq}_{\text{Jensen Inequality}} \mathbb{E}\left[\left(\mathbb{E}_{s_t\sim \mathcal{D}}\left[\left|\left\langle\nabla f\left(\theta^t\right), s_t\right\rangle\right| \mid \theta^t\right]\right)^2\right] \\
&\underbrace{\geq}_{\text{\Cref{Ass5}}} \mu_{\mathcal{D}}^2 \mathbb{E}\left[\left\|\nabla f\left(\theta^t\right)\right\|_{2}^2\right].
\end{align*}




It holds that:
$
\mathbb{E}\left[f\left(\theta^{t+1} \right)\right]- f(\theta^{*}) \leq \mathbb{E}\left[f\left(\theta^{t} \right)- f(\theta^{*})\right]-\frac{\mu_{\mathcal{D}}^2}{2 L} \mathbb{E}\left[\left\|\nabla f\left( \theta^t \right)\right\|_2^2\right]+\frac{L h^{-2t}}{8}.
$

By \Cref{Ass4}, we have $\left\|\nabla f\left(\theta^t\right)\right\|_2^2 \geq 2 \mu\left(f\left(\theta^t\right)-f(\theta^{*}) \right)$, then:

$$
\mathbb{E}\left[f\left(\theta^{t+1} \right)- f(\theta^{*})\right] \leq \left(1-\frac{\mu_{\mathcal{D}}^2 \mu}{L}\right) \mathbb{E}\left[f\left(\theta^{t} \right)- f(\theta^{*})\right]+\frac{L h^{-2t}}{8}.
$$ 
Thus by induction we obtain: $$\forall T \ge 2, \;
\mathbb{E}[f(\theta^{T})-f(\theta^{*})  ]  \leq\left(1-\frac{\mu_{\mathcal{D}}^2 \mu}{L}\right)^{T-1} [f(\theta^{1})-f(\theta^{*})  ]+ \frac{L}{8}\sum_{i=1}^{T-1}\left(1-\frac{\mu_{\mathcal{D}}^2 \mu}{L}\right)^{T-1-i} h^{-2i}.$$
\end{proof}

\begin{proof}[Proof of \Cref{Thm7}]
Since $h \in \left(\frac{1}{\sqrt{1-\frac{\mu_{\mathcal{D}}^2 \mu}{L}}}, \infty \right)$, it holds that: $ \forall T \ge 2, \;  1-\frac{1}{h^2\left(1-\frac{\mu_{\mathcal{D}}^2 \mu}{L}\right)   } >0.$

Then, using \Cref{Lemma6}, we get: 
\begin{align*}
\forall T \ge 2, \; 
\mathbb{E}[f(\theta^{T})-f(\theta^{*})  ]  &\leq\left(1-\frac{\mu_{\mathcal{D}}^2 \mu}{L}\right)^{T-1} [f(\theta^{1})-f(\theta^{*})  ]+\frac{L}{8h^2} \frac{ \left(1-\frac{\mu_{\mathcal{D}}^2 \mu}{L}\right)^{T-2}}{1- \frac{1}{h^2\left(1-\frac{\mu_{\mathcal{D}}^2 \mu}{L}\right)   } }   \\
&\label{eq-strong}\leq \left(1-\frac{\mu_{\mathcal{D}}^2 \mu}{L}\right)^{T-1} \bigg[f(\theta^{1})-f(\theta^{*}) 
+\frac{L}{8} 
%\left(1-\frac{\mu_{\mathcal{D}}^2 \mu}{L}\right)^{T-1} 
\frac{ 1}{h^2 \left(1-\frac{\mu_{\mathcal{D}}^2 \mu}{L} \right)- 1 } \bigg],
\end{align*}
which gives the desired inequality. 

In the particular case where $\mu_{\mathcal{D}} = \frac{K}{\sqrt{d}}$, by replacing $h$ and $\mu_{\mathcal{D}}$ by their formulas, we obtain the desired rate $O\left(\left(1 - \frac{\mu K^2}{dL}\right)^T\right)$.
\end{proof}

%\textbf{Proof of \Cref{Thm6}.}
\begin{proof}[Proof of \Cref{Thm6}]

Let $s \in (0,1)$, we consider $a=1-s\frac{\mu_{\mathcal{D}}^2 \mu }{L}$. By multiplying the inequality~(\ref{eq-strong-exp}) by $a^{-T}$, we get that for all $T \ge 2$:
$$\mathbb{E}[a^{-T} \left( f(\theta^{T})-f(\theta^{*})  \right)]  
\le  \left(a^{-1}-\frac{a^{-1}\mu_{\mathcal{D}}^2 \mu}{L}\right)^{T-1} \bigg[a^{-1}\left(f(\theta^{1})-f(\theta^{*})\right) +\frac{a^{-1} L}{8} \frac{1}{h^2 \left(1-\frac{\mu_{\mathcal{D}}^2 \mu}{L} \right)- 1 }\bigg].$$

As  $\;a^{-1} - \frac{a^{-1} \mu_{\mathcal{D}}^2 \mu }{L} = \frac{1 - \frac{\mu_{\mathcal{D}}^2 \mu}{L}}{1 - s\frac{\mu_{\mathcal{D}}^2 \mu}{L}}  \in (0,1)$
, it holds that: $$\mathbb{E}[\sum_{T=2}^{\infty} a^{-T}\left(f(\theta^{T})-f(\theta^{*})\right) ] = \sum_{T=2}^{\infty}  \mathbb{E}[a^{-T}\left(f(\theta^{T})-f(\theta^{*})\right)] <\infty.$$ 
Therefore: $\sum_{T=2}^{\infty} a^{-T}\left(f(\theta^{T})-f(\theta^{*})\right) <\infty \;\; \text{ a.s.} $, and we conclude that: $$f(\theta^T)-f(\theta^{*}) =o\left(a^{T}\right)\;\; \text{ a.s.} $$
\end{proof}


\end{document}
